%************************************************************************
% Worksheet: GPIO output stages (push-pull, open drain, open collector)
% simulated in Falstad CircuitJS. Circuits: examples/interfaces/falstad/
%
%   make worksheet      builds Worksheet_IO_Ports.pdf (students)
%                       and Worksheet_IO_Ports_Solutions.pdf (lecturer)
%************************************************************************
\documentclass[a4paper,11pt]{article}
%************************************************************************
% Common preamble for the script (docu.tex) and the slides (slides.tex).
% Each root file loads its document class, then %************************************************************************
% Common preamble for the script (docu.tex) and the slides (slides.tex).
% Each root file loads its document class, then %************************************************************************
% Common preamble for the script (docu.tex) and the slides (slides.tex).
% Each root file loads its document class, then \input{preamble},
% then hyperref with its own options, then its page geometry.
%************************************************************************
\usepackage{iftex}
\ifPDFTeX
  % fallback without the RWU font: make PDFTEX=1 (or plain pdflatex)
  \usepackage[utf8]{inputenc}
  \usepackage[T1]{fontenc}
  \usepackage{lmodern}
\else
  % RWU corporate font Barlow Semi Condensed, shipped in fonts/ (SIL OFL)
  \usepackage{fontspec}
  \defaultfontfeatures{Ligatures=TeX}
  \newcommand{\rwufontsetup}[1]{#1{BarlowSemiCondensed}[
      Path=fonts/, Extension=.otf,
      UprightFont=*-Regular, ItalicFont=*-Italic,
      BoldFont=*-SemiBold, BoldItalicFont=*-SemiBoldItalic]}
  \rwufontsetup\setmainfont
  \rwufontsetup\setsansfont
  \setmonofont{Latin Modern Mono}[Ligatures=TeXOff]   % keep -- as two hyphens in code
\fi
\usepackage{xcolor}
\usepackage{color,epsfig,graphicx,amsmath,amsfonts,latexsym,longtable,epic,eepic,fancybox,moreverb,url,version,ifthen,alltt}
\usepackage{shadow}

% RWU corporate colours (as in the RWU LaTeX style, rwudefs.sty)
\definecolor{rwuviolet}{RGB}{102,56,182}
\definecolor{rwuvioletlight}{RGB}{155,125,212}
\definecolor{rwucyan}{RGB}{0,169,206}
\definecolor{rwucyanlight}{RGB}{5,195,222}
\definecolor{rwucyan40}{RGB}{153,221,253}
\colorlet{rwucyandark}{rwucyan!75!black}   % cyan for text on white
\colorlet{rwugray}{black!55}
% pulls the first bullet up under a topic label: \item[{\rot\bf Topic}]~\\[-\lblskip]
% (em-based, so it does not depend on the x-height of the font; was -3ex with Latin Modern)
\newcommand{\lblskip}{1.1em}
% same before a framed listing, which needs room for its top rule
\newcommand{\lstskip}{0.5em}
\newcommand{\rwulogo}[1][height=1.3cm]{\includegraphics[#1]{rwu_logo_hor_lila-cyan_cmyk}}
\usepackage{chngcntr}
\usepackage{verbatim}
\usepackage{tikz}                % own drawings (block diagrams, signals)
\usepackage{adjustbox}           % scale a drawing to fit width and height
\usepackage{tabularx}            % tables with wrapping columns (type L)
\newcolumntype{L}{>{\raggedright\arraybackslash}X}
\usetikzlibrary{arrows.meta,positioning,calc,decorations.pathreplacing,shapes.symbols,shapes.geometric}
\usepackage{tikz-timing}         % timing diagrams (UART, I2C, SPI, ...)

\usepackage{listings}
\colorlet{codegreen}{black!55}
\definecolor{codegray}{rgb}{0.5,0.5,0.5}
\colorlet{codepurple}{rwucyandark}
\definecolor{backcolour}{rgb}{0.95,0.95,0.92}
\lstdefinestyle{mystyle}{
    backgroundcolor=\color{backcolour},
    commentstyle=\color{codegreen},
    keywordstyle=\color{rwuviolet}\bfseries,
    numberstyle=\tiny\color{codegray},
    stringstyle=\color{codepurple},
    basicstyle=\ttfamily\footnotesize,
    breakatwhitespace=false,
    breaklines=true,
    captionpos=b,
    keepspaces=true,
    numbers=left,
    numbersep=5pt,
    showspaces=false,
    showstringspaces=false,
    showtabs=false,
    tabsize=2
}
\lstset{style=mystyle}

\colorlet{mygreen}{black!55}
\definecolor{mygray}{rgb}{0.5,0.5,0.5}
\colorlet{mymauve}{rwucyandark}

%Customize a bit the look
\lstset{ %
backgroundcolor=\color{white}, % choose the background color; you must add \usepackage{color} or \usepackage{xcolor}
basicstyle=\footnotesize, % the size of the fonts that are used for the code
breakatwhitespace=false, % sets if automatic breaks should only happen at whitespace
breaklines=true, % sets automatic line breaking
captionpos=b, % sets the caption-position to bottom
commentstyle=\color{mygreen}, % comment style
deletekeywords={...}, % if you want to delete keywords from the given language
escapeinside={\%*}{*)}, % if you want to add LaTeX within your code
extendedchars=true, % lets you use non-ASCII characters; for 8-bits encodings only, does not work with UTF-8
frame=single, % adds a frame around the code
keepspaces=true, % keeps spaces in text, useful for keeping indentation of code (possibly needs columns=flexible)
keywordstyle=\color{rwuviolet}\bfseries, % keyword style
% language=Octave, % the language of the code
morekeywords={*,...}, % if you want to add more keywords to the set
numbers=left, % where to put the line-numbers; possible values are (none, left, right)
numbersep=5pt, % how far the line-numbers are from the code
numberstyle=\tiny\color{mygray}, % the style that is used for the line-numbers
rulecolor=\color{rwucyan40}, % if not set, the frame-color may be changed on line-breaks within not-black text (e.g. comments (green here))
showspaces=false, % show spaces everywhere adding particular underscores; it overrides 'showstringspaces'
showstringspaces=false, % underline spaces within strings only
showtabs=false, % show tabs within strings adding particular underscores
stepnumber=1, % the step between two line-numbers. If it's 1, each line will be numbered
stringstyle=\color{mymauve}, % string literal style
tabsize=2, % sets default tabsize to 2 spaces
title=\lstname % show the filename of files included with \lstinputlisting; also try caption instead of title
}
%END of listing package%

\definecolor{darkgray}{rgb}{.4,.4,.4}
\definecolor{purple}{rgb}{0.65, 0.12, 0.82}

%define Javascript language
\lstdefinelanguage{JavaScript}{
keywords={typeof, new, true, false, catch, function, return, null, catch, switch, var, if, in, while, do, else, case, break},
keywordstyle=\color{blue}\bfseries,
ndkeywords={class, export, boolean, throw, implements, import, this},
ndkeywordstyle=\color{darkgray}\bfseries,
identifierstyle=\color{black},
sensitive=false,
comment=[l]{//},
morecomment=[s]{/*}{*/},
commentstyle=\color{purple}\ttfamily,
stringstyle=\color{red}\ttfamily,
morestring=[b]',
morestring=[b]"
}

\lstset{
language=JavaScript,
extendedchars=true,
basicstyle=\footnotesize\ttfamily,
showstringspaces=false,
showspaces=false,
numbers=left,
numberstyle=\footnotesize,
numbersep=9pt,
tabsize=2,
breaklines=true,
showtabs=false,
captionpos=b,
columns=fullflexible   % natural widths, also for spaces (comments are set smaller)
}

%define Rust language (listings has no built-in Rust)
\lstdefinelanguage{Rust}{
keywords={as, break, const, continue, crate, else, enum, extern, false, fn, for, if, impl, in, let, loop, match, mod, move, mut, pub, ref, return, self, Self, static, struct, super, trait, true, type, unsafe, use, where, while, async, await, dyn},
sensitive=true,
comment=[l]{//},
morecomment=[s]{/*}{*/},
morestring=[b]"
}
,
% then hyperref with its own options, then its page geometry.
%************************************************************************
\usepackage{iftex}
\ifPDFTeX
  % fallback without the RWU font: make PDFTEX=1 (or plain pdflatex)
  \usepackage[utf8]{inputenc}
  \usepackage[T1]{fontenc}
  \usepackage{lmodern}
\else
  % RWU corporate font Barlow Semi Condensed, shipped in fonts/ (SIL OFL)
  \usepackage{fontspec}
  \defaultfontfeatures{Ligatures=TeX}
  \newcommand{\rwufontsetup}[1]{#1{BarlowSemiCondensed}[
      Path=fonts/, Extension=.otf,
      UprightFont=*-Regular, ItalicFont=*-Italic,
      BoldFont=*-SemiBold, BoldItalicFont=*-SemiBoldItalic]}
  \rwufontsetup\setmainfont
  \rwufontsetup\setsansfont
  \setmonofont{Latin Modern Mono}[Ligatures=TeXOff]   % keep -- as two hyphens in code
\fi
\usepackage{xcolor}
\usepackage{color,epsfig,graphicx,amsmath,amsfonts,latexsym,longtable,epic,eepic,fancybox,moreverb,url,version,ifthen,alltt}
\usepackage{shadow}

% RWU corporate colours (as in the RWU LaTeX style, rwudefs.sty)
\definecolor{rwuviolet}{RGB}{102,56,182}
\definecolor{rwuvioletlight}{RGB}{155,125,212}
\definecolor{rwucyan}{RGB}{0,169,206}
\definecolor{rwucyanlight}{RGB}{5,195,222}
\definecolor{rwucyan40}{RGB}{153,221,253}
\colorlet{rwucyandark}{rwucyan!75!black}   % cyan for text on white
\colorlet{rwugray}{black!55}
% pulls the first bullet up under a topic label: \item[{\rot\bf Topic}]~\\[-\lblskip]
% (em-based, so it does not depend on the x-height of the font; was -3ex with Latin Modern)
\newcommand{\lblskip}{1.1em}
% same before a framed listing, which needs room for its top rule
\newcommand{\lstskip}{0.5em}
\newcommand{\rwulogo}[1][height=1.3cm]{\includegraphics[#1]{rwu_logo_hor_lila-cyan_cmyk}}
\usepackage{chngcntr}
\usepackage{verbatim}
\usepackage{tikz}                % own drawings (block diagrams, signals)
\usepackage{adjustbox}           % scale a drawing to fit width and height
\usepackage{tabularx}            % tables with wrapping columns (type L)
\newcolumntype{L}{>{\raggedright\arraybackslash}X}
\usetikzlibrary{arrows.meta,positioning,calc,decorations.pathreplacing,shapes.symbols,shapes.geometric}
\usepackage{tikz-timing}         % timing diagrams (UART, I2C, SPI, ...)

\usepackage{listings}
\colorlet{codegreen}{black!55}
\definecolor{codegray}{rgb}{0.5,0.5,0.5}
\colorlet{codepurple}{rwucyandark}
\definecolor{backcolour}{rgb}{0.95,0.95,0.92}
\lstdefinestyle{mystyle}{
    backgroundcolor=\color{backcolour},
    commentstyle=\color{codegreen},
    keywordstyle=\color{rwuviolet}\bfseries,
    numberstyle=\tiny\color{codegray},
    stringstyle=\color{codepurple},
    basicstyle=\ttfamily\footnotesize,
    breakatwhitespace=false,
    breaklines=true,
    captionpos=b,
    keepspaces=true,
    numbers=left,
    numbersep=5pt,
    showspaces=false,
    showstringspaces=false,
    showtabs=false,
    tabsize=2
}
\lstset{style=mystyle}

\colorlet{mygreen}{black!55}
\definecolor{mygray}{rgb}{0.5,0.5,0.5}
\colorlet{mymauve}{rwucyandark}

%Customize a bit the look
\lstset{ %
backgroundcolor=\color{white}, % choose the background color; you must add \usepackage{color} or \usepackage{xcolor}
basicstyle=\footnotesize, % the size of the fonts that are used for the code
breakatwhitespace=false, % sets if automatic breaks should only happen at whitespace
breaklines=true, % sets automatic line breaking
captionpos=b, % sets the caption-position to bottom
commentstyle=\color{mygreen}, % comment style
deletekeywords={...}, % if you want to delete keywords from the given language
escapeinside={\%*}{*)}, % if you want to add LaTeX within your code
extendedchars=true, % lets you use non-ASCII characters; for 8-bits encodings only, does not work with UTF-8
frame=single, % adds a frame around the code
keepspaces=true, % keeps spaces in text, useful for keeping indentation of code (possibly needs columns=flexible)
keywordstyle=\color{rwuviolet}\bfseries, % keyword style
% language=Octave, % the language of the code
morekeywords={*,...}, % if you want to add more keywords to the set
numbers=left, % where to put the line-numbers; possible values are (none, left, right)
numbersep=5pt, % how far the line-numbers are from the code
numberstyle=\tiny\color{mygray}, % the style that is used for the line-numbers
rulecolor=\color{rwucyan40}, % if not set, the frame-color may be changed on line-breaks within not-black text (e.g. comments (green here))
showspaces=false, % show spaces everywhere adding particular underscores; it overrides 'showstringspaces'
showstringspaces=false, % underline spaces within strings only
showtabs=false, % show tabs within strings adding particular underscores
stepnumber=1, % the step between two line-numbers. If it's 1, each line will be numbered
stringstyle=\color{mymauve}, % string literal style
tabsize=2, % sets default tabsize to 2 spaces
title=\lstname % show the filename of files included with \lstinputlisting; also try caption instead of title
}
%END of listing package%

\definecolor{darkgray}{rgb}{.4,.4,.4}
\definecolor{purple}{rgb}{0.65, 0.12, 0.82}

%define Javascript language
\lstdefinelanguage{JavaScript}{
keywords={typeof, new, true, false, catch, function, return, null, catch, switch, var, if, in, while, do, else, case, break},
keywordstyle=\color{blue}\bfseries,
ndkeywords={class, export, boolean, throw, implements, import, this},
ndkeywordstyle=\color{darkgray}\bfseries,
identifierstyle=\color{black},
sensitive=false,
comment=[l]{//},
morecomment=[s]{/*}{*/},
commentstyle=\color{purple}\ttfamily,
stringstyle=\color{red}\ttfamily,
morestring=[b]',
morestring=[b]"
}

\lstset{
language=JavaScript,
extendedchars=true,
basicstyle=\footnotesize\ttfamily,
showstringspaces=false,
showspaces=false,
numbers=left,
numberstyle=\footnotesize,
numbersep=9pt,
tabsize=2,
breaklines=true,
showtabs=false,
captionpos=b,
columns=fullflexible   % natural widths, also for spaces (comments are set smaller)
}

%define Rust language (listings has no built-in Rust)
\lstdefinelanguage{Rust}{
keywords={as, break, const, continue, crate, else, enum, extern, false, fn, for, if, impl, in, let, loop, match, mod, move, mut, pub, ref, return, self, Self, static, struct, super, trait, true, type, unsafe, use, where, while, async, await, dyn},
sensitive=true,
comment=[l]{//},
morecomment=[s]{/*}{*/},
morestring=[b]"
}
,
% then hyperref with its own options, then its page geometry.
%************************************************************************
\usepackage{iftex}
\ifPDFTeX
  % fallback without the RWU font: make PDFTEX=1 (or plain pdflatex)
  \usepackage[utf8]{inputenc}
  \usepackage[T1]{fontenc}
  \usepackage{lmodern}
\else
  % RWU corporate font Barlow Semi Condensed, shipped in fonts/ (SIL OFL)
  \usepackage{fontspec}
  \defaultfontfeatures{Ligatures=TeX}
  \newcommand{\rwufontsetup}[1]{#1{BarlowSemiCondensed}[
      Path=fonts/, Extension=.otf,
      UprightFont=*-Regular, ItalicFont=*-Italic,
      BoldFont=*-SemiBold, BoldItalicFont=*-SemiBoldItalic]}
  \rwufontsetup\setmainfont
  \rwufontsetup\setsansfont
  \setmonofont{Latin Modern Mono}[Ligatures=TeXOff]   % keep -- as two hyphens in code
\fi
\usepackage{xcolor}
\usepackage{color,epsfig,graphicx,amsmath,amsfonts,latexsym,longtable,epic,eepic,fancybox,moreverb,url,version,ifthen,alltt}
\usepackage{shadow}

% RWU corporate colours (as in the RWU LaTeX style, rwudefs.sty)
\definecolor{rwuviolet}{RGB}{102,56,182}
\definecolor{rwuvioletlight}{RGB}{155,125,212}
\definecolor{rwucyan}{RGB}{0,169,206}
\definecolor{rwucyanlight}{RGB}{5,195,222}
\definecolor{rwucyan40}{RGB}{153,221,253}
\colorlet{rwucyandark}{rwucyan!75!black}   % cyan for text on white
\colorlet{rwugray}{black!55}
% pulls the first bullet up under a topic label: \item[{\rot\bf Topic}]~\\[-\lblskip]
% (em-based, so it does not depend on the x-height of the font; was -3ex with Latin Modern)
\newcommand{\lblskip}{1.1em}
% same before a framed listing, which needs room for its top rule
\newcommand{\lstskip}{0.5em}
\newcommand{\rwulogo}[1][height=1.3cm]{\includegraphics[#1]{rwu_logo_hor_lila-cyan_cmyk}}
\usepackage{chngcntr}
\usepackage{verbatim}
\usepackage{tikz}                % own drawings (block diagrams, signals)
\usepackage{adjustbox}           % scale a drawing to fit width and height
\usepackage{tabularx}            % tables with wrapping columns (type L)
\newcolumntype{L}{>{\raggedright\arraybackslash}X}
\usetikzlibrary{arrows.meta,positioning,calc,decorations.pathreplacing,shapes.symbols,shapes.geometric}
\usepackage{tikz-timing}         % timing diagrams (UART, I2C, SPI, ...)

\usepackage{listings}
\colorlet{codegreen}{black!55}
\definecolor{codegray}{rgb}{0.5,0.5,0.5}
\colorlet{codepurple}{rwucyandark}
\definecolor{backcolour}{rgb}{0.95,0.95,0.92}
\lstdefinestyle{mystyle}{
    backgroundcolor=\color{backcolour},
    commentstyle=\color{codegreen},
    keywordstyle=\color{rwuviolet}\bfseries,
    numberstyle=\tiny\color{codegray},
    stringstyle=\color{codepurple},
    basicstyle=\ttfamily\footnotesize,
    breakatwhitespace=false,
    breaklines=true,
    captionpos=b,
    keepspaces=true,
    numbers=left,
    numbersep=5pt,
    showspaces=false,
    showstringspaces=false,
    showtabs=false,
    tabsize=2
}
\lstset{style=mystyle}

\colorlet{mygreen}{black!55}
\definecolor{mygray}{rgb}{0.5,0.5,0.5}
\colorlet{mymauve}{rwucyandark}

%Customize a bit the look
\lstset{ %
backgroundcolor=\color{white}, % choose the background color; you must add \usepackage{color} or \usepackage{xcolor}
basicstyle=\footnotesize, % the size of the fonts that are used for the code
breakatwhitespace=false, % sets if automatic breaks should only happen at whitespace
breaklines=true, % sets automatic line breaking
captionpos=b, % sets the caption-position to bottom
commentstyle=\color{mygreen}, % comment style
deletekeywords={...}, % if you want to delete keywords from the given language
escapeinside={\%*}{*)}, % if you want to add LaTeX within your code
extendedchars=true, % lets you use non-ASCII characters; for 8-bits encodings only, does not work with UTF-8
frame=single, % adds a frame around the code
keepspaces=true, % keeps spaces in text, useful for keeping indentation of code (possibly needs columns=flexible)
keywordstyle=\color{rwuviolet}\bfseries, % keyword style
% language=Octave, % the language of the code
morekeywords={*,...}, % if you want to add more keywords to the set
numbers=left, % where to put the line-numbers; possible values are (none, left, right)
numbersep=5pt, % how far the line-numbers are from the code
numberstyle=\tiny\color{mygray}, % the style that is used for the line-numbers
rulecolor=\color{rwucyan40}, % if not set, the frame-color may be changed on line-breaks within not-black text (e.g. comments (green here))
showspaces=false, % show spaces everywhere adding particular underscores; it overrides 'showstringspaces'
showstringspaces=false, % underline spaces within strings only
showtabs=false, % show tabs within strings adding particular underscores
stepnumber=1, % the step between two line-numbers. If it's 1, each line will be numbered
stringstyle=\color{mymauve}, % string literal style
tabsize=2, % sets default tabsize to 2 spaces
title=\lstname % show the filename of files included with \lstinputlisting; also try caption instead of title
}
%END of listing package%

\definecolor{darkgray}{rgb}{.4,.4,.4}
\definecolor{purple}{rgb}{0.65, 0.12, 0.82}

%define Javascript language
\lstdefinelanguage{JavaScript}{
keywords={typeof, new, true, false, catch, function, return, null, catch, switch, var, if, in, while, do, else, case, break},
keywordstyle=\color{blue}\bfseries,
ndkeywords={class, export, boolean, throw, implements, import, this},
ndkeywordstyle=\color{darkgray}\bfseries,
identifierstyle=\color{black},
sensitive=false,
comment=[l]{//},
morecomment=[s]{/*}{*/},
commentstyle=\color{purple}\ttfamily,
stringstyle=\color{red}\ttfamily,
morestring=[b]',
morestring=[b]"
}

\lstset{
language=JavaScript,
extendedchars=true,
basicstyle=\footnotesize\ttfamily,
showstringspaces=false,
showspaces=false,
numbers=left,
numberstyle=\footnotesize,
numbersep=9pt,
tabsize=2,
breaklines=true,
showtabs=false,
captionpos=b,
columns=fullflexible   % natural widths, also for spaces (comments are set smaller)
}

%define Rust language (listings has no built-in Rust)
\lstdefinelanguage{Rust}{
keywords={as, break, const, continue, crate, else, enum, extern, false, fn, for, if, impl, in, let, loop, match, mod, move, mut, pub, ref, return, self, Self, static, struct, super, trait, true, type, unsafe, use, where, while, async, await, dyn},
sensitive=true,
comment=[l]{//},
morecomment=[s]{/*}{*/},
morestring=[b]"
}

\usepackage[colorlinks=true,pdfpagemode=UseNone]{hyperref}
\hypersetup{linkcolor=rwuviolet,citecolor=rwucyandark,urlcolor=rwucyandark}
\usepackage[a4paper,top=22mm,bottom=25mm,left=22mm,right=22mm]{geometry}
\usepackage{titlesec}
\titleformat*{\section}{\sffamily\Large\bfseries\color{rwuviolet}}
\titleformat*{\subsection}{\sffamily\large\bfseries\color{rwucyandark}}
\usepackage{fancyhdr}
\pagestyle{fancy}
\fancyhf{}
\fancyhead[L]{\color{rwugray}\small Embedded Computing -- Worksheet: GPIO output stages\ifdefined\solutions\ (solutions)\fi}
\fancyhead[R]{\color{rwugray}\small RWU}
\fancyfoot[C]{\color{rwugray}\small\thepage}
\renewcommand{\headrulewidth}{0pt}
\usepackage{booktabs,enumitem,siunitx,amssymb,needspace}
\sisetup{per-mode=symbol}
\usepackage[most]{tcolorbox}
\setlength{\parindent}{0pt}
\setlength{\parskip}{0.5ex}
\setlist{itemsep=0.3ex,topsep=0.5ex}

% generated by examples/interfaces/falstad/make_circuits.py
\def\linkA{https://www.falstad.com/circuit/circuitjs.html?ctz=DwYwlgTgBAZgvAIgIwKgFwM6IEwFMC0ADAOypgiJKEB02qIARogKyGoAOTCrqAbhC1QBbTJQIkEAPgBQUKMCFCoAOyGIA5uzAB7fKu0ZU8BGyi80iGgDZUDXJeqEkdKBi6mAJhgCel2x8pbCxNUNwCQgHoZOQUlVQ0tXXYhAyM-M2DrW3sTR0IADlD3VC9fEKgGcJQKzKLwtijZeQAPKGbEABZ8qABmbChCAbTyjDAcUzQcgEEALigABQBXDAALfHZFgBtNqABxeYBJAHkobUW0DbQoEBWAQwhbkEmIMAw0cigAChBFiF5cdAPEC4CAASgQjRirXaCC6UCsHQGQ1g6VGlA6qEmiAAMtp1B8wMpLlA4KdzsSILh8W8QRUwGhqFAACorAEYADuuFw7FcZwgwNg2n5uAwUAAap9DgA5UGwCDaJREajMcXobS9ajdMXUCHRFptTrdfLdQamYymNHIDHoHLY3AwK4YEDadi4OYADRJ4slBxlABooABNL1aZTXX6U5RXT4AMnUaAA3AM5mcrtoYOhWVBQ6DGQAlMDqFaO52uubsbSc6CEhYAWSOAGUoLdlB4oFL6w2dZD5NiDchsN1sMwrFB8oNhzYUeV2IEoCsxggetQeqhNto-D3gAd+84hyOoEhiKPJ8MLZsHMxUAvEMvV1uMzDsB0JwVeoREdgCsNqikEjp1hSQwH37Z9Bh6D930RCDrXNYQNwQTQAP0YC9WAdlQJfKAvyHLDT2nBo0Iwp88IPMCoLPXUmnQ-sYOw58oOwkdKK3Yjb0gydGMHQoCKomJ1FoyC+nAoTmBcOCtzzQTESQABOfo6KQKwzXSdk4KgIRbmaXhbxXPj5DY5BjyYk9SKnCSiOkkzeisfdzM3NDNjDGEels6zXKHZ9hmtLEEGlfTgCmRRcGeWi3M4joERM4ZVw0nILQoSI0N4ftIug-JETS6LeKgNTbyMAQEGHUxNO0nBqCQK86VuLBkGoKwry3ASYSynoMqgLKOiUli0PVXBlEoUhXFnZAf1k60RrYPUYh5HSRkSjEexmqAcm6sx8tcRKaEIRq0IibRepWgbuFCEaeOMFwRpQab5FmxAhqdBwUCW26VsQKc5oeraKoC-aZGACJwAgGQgA}
\def\linkB{https://www.falstad.com/circuit/circuitjs.html?ctz=DwYwlgTgBAZgvAIgIwKgFwM6KQUwLRIAMqYI2hAdAEyogBGiArMVAA4MLOoBuETqAW0xN8SFAD4AUFCjABAqADsBiAOaswAezzLNGVPAQtuaRJQBsqOjjMVCSGlAwcWAEwwBPM1dfYrpo1RnX0CAeikZOQVlNQ1tVgE9A28oE1tLKGtbQkIADiCXVHcvQMyQlEyAlmDvcOlZAA8oBsQAFlyoAGYqKEJe5NKMMEQqFjQbBAAhAC4oAAUAVwwACzxWBYAbDahNBbR1tCcAdzA0EGWwRVUoAEMoEBvWG-A0MG4cKA3NG9cEOsimi0EO0oOZWr1+rAUkNsK1UONEAAVZYfDAARwWNwgHyON3eUDAGHQKLYEE0qggNwUaHJqg2l2uaBJGkUAG5iR9WNjXBA3jhoBgNjgjlBKeMnDhMFBlpoRTAbhhDkyPqobuMiRgTmdlhQ-hFGs02h1ch0+ixDNVhsg4egJgAZHAwQ4YECaVg4Waq8UAcVuilc8wAkgA5ChQABKYFUy2drvdswwC1YrA2HnuCwg2MUhxu-oJUrx-JuqhwYYAyvTXPyiZpFBzRVGY7r-rJw4aEABOcxdUZQEHdc0pI6GRwwPjIHKTygsAQ3BrcbAUcyMKxgBWL5d6+rAEVArs9voODoDgbEFvAQPto8HqBIADs3ZPUMGGxGqAuiE6FE6W8iu+wD43qMx6jKev6yBsdZAsBQF5FAVAmgMNoIgg3rgcAMDtlQrR9FQd49J0hDgnho5+FAiSxFoayJPo56YdBOFdOYx5EUx+TPjOmiUdoui0fqO5YYxJHwYxT4Wuh-4INhfQDiJMnMWB56SYRxEml0rFibU-HKRpvYqV0jCkWE-HXEC+mdLk4L6a0SAZOJSlYXBsnSQein8W2Zmsfej5ea0NriVAw4pLO86ft+6EAILyJK-Ltvp3nqcRjB2Z+ggTNUZDGduOlWXpzGuRxElxflsmtN0BX2fxkHFSB8kdAh7GGMhExzCG6EgO2ZUEb2XV9jkYH3NgohniZnXlThfS9V1bnbj8ABWUATFQFQ4FaLA4IoiAAMKPM8py5iAHwABQAGIAJSCJtyCiI4AhArgeCEB2QQBHa3wBg8TwvAdHxbQA+naqBfKofjnvNi2pYta2oBtiAndiGIbSAaZHQAEgAXhd5FXUQU52IIQLMFOOT404ASIrSQqwAjCxIx4QPkqD-HgxMK5QyksMIBWwoRmqx0AGqhIoGBYwIV3ThUd38GTiBzNyvLvAKQoimKAtCyLDMg8g6GaItV13kErCfkZjhG4E+qRKwqTQplcItpbENScYNttOhoSaPxuuc-ehvGwMSCdEQd4-mwfgW7IVsLggIcutkK7nu7UjAKE4AQFIQA}
\def\linkC{https://www.falstad.com/circuit/circuitjs.html?ctz=DwYwlgTgBAZgvAIgIwKgFwM6KQUwLQAMAnKmCNgQHQBMqIARogKwGoAOjCLqAbhM6gC2mbPmIIAfACgoUYIMFQAdoMQBzNmAD2eFVoyp4CVlB5pEVAGyp6OC5QJJaUDJxMATDAE8LN99htzY1RXf2CAemlZeUUVdU0dNkF9Q19TIKsbO2MHAgAOELdUTx9gqHowlHKMwrDWSJk5AA8oJsQAFjyoAGZqKAJ+1LKMMERqEzRsgGEALigAFQB3LSg2AFcMAAs8dYAbXagtNbR1zEOlc5woXbAlK4AKeg2oEC0lSfftJQBKBAbolptBCdKCWdr9QawNIjbDtVCTRAAIS0aE2q1uGCgAEMIFdFpAcO50Fo1DhUTgIAAaKA4ShqSjoZZQdw4HhkHCYt7YlybHGE663K5aaBY9EXLbCyZEtArUV3MBqTb0I4QSh-KLNVodLp5LoDExGEww5Bw9DZADKZKgAEEoHAoAAJbFKImIu1QAAyUHuwtM7NMFIwWO+DPNrzYODmNzuLzWEFx72pbC0iwpUFu6E2VwACgBZADy5sOMBtzulWagADkC0WtCXEWr-nIvUCnF1qExLFA8gMO9YoWU2AEoJtRghupRuqhdlpfE3gABJLXIajtztQJAAdi7faGRt29iYqFHiAnU-nJaB1HavfyPQI4PGBQHVWS8W0O2SBgvy+vA26D73uCAGmoaQizggGgfno34asAiy-jeUBPshSG7gO9RwQhV5oeuf5AXu6qNPBy4gch15AchnaEfO2GnoBfaUauz5gfOaikYBvT-pxTDOKxcEAEoceCSBEH0ZFIJYBppIsYFQIIWJNDwp6TkR0R0cg25UTuuH9vxxEtswq7aVAnS9tRGHsGkJ7jqp1wQZhxFLkCvFrl2TBaehckYAeORHiOY5nmpciXogYK3l0m4MfkQyvhBUGJF+wXAKFCDhZRUXAQ+hHyfFCS6EltHLulKElRZ+nqcVumgkhIE0Vhy6ZeR4JNV5c4NUCrXrq1ur1cR7GdZx4xQE13S8X10RCYNIliSNgGSdJZSyWkClKSp54dcwnnrmVentcRGlkYx3TdDu5X7dEuwXECJ1nV2t3NUMpoIgg2YLpWAD61rJYdp0mQ97SSRNcjWgoZJpjdf2A+5xnQ0MU7ydkRrkBEm1cLDkkjVDQOWUVnXY12m5-W1qMHY1xPdQx52k5d12IETd1Yzu15PfC2RvZ9iLJSsOBKGMz4YEOyB7cgVnBBq0RsKYp4hCjVAEEe87hFocE83zXAhELONGM4WuSI0kvSwgm6y-YKBNob2Rbrw9Omzk5twcr0jAOE4AQNIQA}
\def\linkD{https://www.falstad.com/circuit/circuitjs.html?ctz=DwYwlgTgBAZgvAIgIwKgFwM6IEwFMC0ADAJypgiJKEB02qIARogKyGoAOTCrqAbhC1QBbTJQIkEAPgBQUKMCFCoAOyGIA5uzAB7fKu0ZU8BGyi80iGgDZUDXJeqEkdKBi6mAJhgCel2x8pbCxNUNwCQgHoZOQUlVQ0tXXYhAyM-M2DrW3sTR0IADlD3VC9fEKgGcJQKzKLwtijZeQAPKGbEABZ8qABmbChCAbTyjDAcUzQcgBEALigAeXZcZXwPCABDMGUKgFcMKAAKAElsAGF8ABswAGtcAEozXAgMPah2PYALfHeLi4RGmKtdoILpQKwdAZDWDpUaUDqoSaIKa4XjkXBQJCuZZeKDrKAgC7aEDXAA0UA8KLRUH6IHW2w+2guHnQH3RDFehIA7odOZBcMyAIIAOSmd2oUAACjtfvgduwoAAlAD6EtxymZ7P2tPY63AaDpIHRpyVmtxEHRGCuFOe1H+0RabU63Xy3UGpmMplhyHh6ByABVWW9Pt9pRdXPr1GzcFzyRAwLxcPs0IGMOshOjtbqwPrlIbYNpoCBtEIdXGMNplOKAMpFpYYuYAIQAqlW1cyJRLq7X0dg5iAdhBzco0FAtoqVbaAfIFY6EMQrNSXVBQdgXcNTJzjC4YAJkIR9-vHMJ1s1eJRqFZmLYwOssMgL1ep8BucD54vus5uqvCtDIvbgEcs6fu+GIAOwLt+66hBcDhXlAHxjAgPTUD0dpNM+QHgSB2AdIMkG-g0-4wLOOF4UuPSEBC+EesI2gJDoegpIYT7qLOFEQuxvSUb0PQ2ARaGAiROG9FYrpQjRriIUgPqIggFKonmKBPgAMrOuFkd0PTMP01HpOwgTwYhyGoVAhJ+E+gHAupIEdCua78a4MG5HBCGIMZAnyMRVl2d0zDOO+UFQCk9G6PozH-qxwJ+f0nHRTxfE0U+QKdH0Ilie6MJSTJOTyVSdBPjOVn7hitnLsVSDOIFm7pEIJ5nkhKEecA0BFYM0kQtZ2DMAl6S7h0oH7k1L5udx2CgTFo3dVBT7DUhk0Lp1U38TNanFV1C5xet03-rNm3dVAzBYVty07bOh0Qft4J4UtiWnVF-ljf0e09X+6EXNsUVHft53UjhwzZYgzZVk1ICzld1KXbh8WBRQCB4PgVBNZFiDgz0vFglDvE-rd6GzW+35QG+0Xbbjs5E-5wHEyd6GWZQq4HRTWFUxJlqwagrkNahT5eTgUMdOCXEdeCwzVMFCCaAxySpNzJFQ91gyxVYGXlGLEuhUxQ2y4M-NUXLlXUzEs2kQzOly0rJOG2x3FdAr3HM+Zd0jRCcWxbZFvyMjc3O1hnFK9jDvoYVTvLqCnEdEgL0bhJtWnm5jUrcCSBMw9evbgHlvArF-k-fbr0Zywyf9ODueEW9H0Fxt2dYTb-0IjkHYg2DqcY4MzDEGn5Sw-DiMsU3rfty3YJ-QbDrAtgwng26gVetJdeIO8GBfD8YbmjATzLIaTXrB4ABWUA5MPuCIaYyyIKc6w6nqBrogcABidzCMoYgI6LifiKQ4aA68mZX7mRomnsVAhJ1CBCfNvPeORw6oCPukU+CBz6X2zNfQ499H7PwqsIN+RAP4YGCEA7QIDkChDvKXGI4D96UD4jA8ocCFSJjALg5BBxtAfCEBgB+QUn5w2wKYIQVkBqDU-ggKUMo5TjglPgwhSl-zkJyCZahJ8uG33NAARx2BvbwhwAASAAvDhQguFUAPEeIKY8DzGM9MEU4hJiSwFUeo3M3hJGgP-NofeXDsAmQwPpZAL0XA+LYPaGI8p6qelhvCKcwSKEglMKE0I4SmoRG0K49xlDQg+IjiLDggQgnyBCW5eJDgnCJOSdIYAERwAQBkEAA}
\def\linkE{https://www.falstad.com/circuit/circuitjs.html?ctz=DwYwlgTgBAZgvAIgIwKgFwM6KQUwLQAMA7KmCNgQHQBMqIARogKwGoAOjCLqAbhM6gC2mbPmIIAfACgoUYIMFQAdoMQBzNmAD2eFVoyp4CVlB5pEVAGyp6OC5QJJaUDJxMATDAE8LN99htzY1RXf2CAemlZeUUVdU0dNkF9Q19TIKsbO2MHAgAOELdUTx9gqHowlHKMwrDWSJk5d1jVBDB3HABDABs8CBw1VLKwLGR8JAAWVAhRkyV7AiYbHl8G6IAPKHXECbyoAGZqKAJjoZMMMERqEzRsgFEALigAeTYcJSh3CE6wD4B3MBoAAWUE6UCYUAAalA2ABXbq9WFsJ6CTrrMCCWGKTQfHhabpoTpqHAINZyTbbBC7KCWCbHU6wNIXbBTdDZABCWmBUAAcgBZZ4AZVB-RcALQICBOHc5S86ClLk6ghwLjAaiUPUoUAACgBJHlIKAjUHKLQQVHdA6UfZQqAAcT1zyNHy0byUeC+Pw+yQ6ABpSVFyVsdns8nsTiYjOdLshWbdEHqeUcwRDIXg0PicN8lGgoK73h7vr8Yb8tQAVBX9NTwzoZ6AgToupTdOUYLSwiAgFUgDv9HNQAAU7m0HQAlJQA41gDxgwgAJyWA5MI4LqBIOfOKOoP5bqCo9YrBD7a2T6JqWer-bLqCXvKsrdk4Ayymr9dHJBERdvs5CLSIdpdL0VankGlLUhMSCLhGP6qiyqDxke1q2lWNZ1iBwB-LOH5fhuUDXHs36MhEgYYbO+Frrh1BEO+G4-o+mGUlRNFHPs+w4ZuqwkQxiCsexUATAQJyEQ+JHdB8lLYRR76fmuliRjs8HZJCAAiynoSAZH5FJeFaUxMHkGMhDWI+56Mbp1E6Xs+wEPenFTtAjEWW+eEWXpREmPwR7WYJ6GmVcrkBSxNl0SRFJXNQdK8fS8llMysaKYgfIAMIAKpQN0WidO46EAIIKDgtwOTsgl4Uwi4CUJtFEfsQjZOcBn1CRMCzlFTH8SVblGFUyTxNoeBJCk9EtWxLksSNwl2dEzWUlFbF7BVByWAU7m-r1Oh6AYJmzgt+x3u1JwQdYK1DWBHVhvtBzXCFU7cVSJWHAd93Ltd0S3Tt1z8bSl0xY1U5idtX0PZ9kXPURcbZImKCPgASgDdIRQdX34TBO5pPuh5LI+RVUkjAnA99MGeRMRA+TDZF47slVHJTKNGJunlIIJgkOEIaKHkglCWEs5RgJ0owc1z6G6uTB15Cc1DnTTK0uN09jc0CMbHjVJ1XJLYsHCNUsiTdw3lers1LS9ci3VFlPgrhWuTcbs5MBb6u2yxhvHU1Nu4XNNJlYty27j1CAaH1G2+bOXOLrtdIh-xkFG6RlIR1eRxx1dzs67HntA9RJwPdH-2Uhn31QHn8dDODCb6rQZOp4uCMe1X+S02jbMCFjwee9Xic-dMOwk79GwtTZ4ITHS0HS3FkwJQgShmha9qOr6eZuoWXp7loHToWFXAyUwg-RTBo8l1wULppm2a5vm7qesWOLoVlABWUDZFRqA4DGJjvIg0M4BcGCEkoXaDloQJBAYFHEIeYVISZCDAt3c4QRtTwkRGwDA6AtDgihKgDKagAiPlvvfYqT8X5PzAR-L+P8-4DgAUAkBe4wHExMIIKBTMQhBHQVoTByAQizGvu4O+2Qx73wIffIhn8RikJVOQwBwDQEUB8nuCSTMmYsxcEEZKaUMpZUHErW04QoC6gAPqqSoRgrBJEUFvwQNQbmGA2CIFpEMZw1jgiBmiGwUwTIDKYycXIFx2Rdq8DcU3Ei4QtAmMEdgKoVibG2XMewAInjgAuMPPVAJU4gkhLMR+EIDjbFESqA42gcSEk8RCA1XIUNAnBKkMAcI4AIDSCAA}
\def\linkF{https://www.falstad.com/circuit/circuitjs.html?ctz=DwYwlgTgBAZgvAIgIwKgFwM6IEwFMC0ADAOypgiJKEB02qIARogKyGoAOTCrqAbhC1QBbTJQIkEAPgBQUKMAAmQqADshiAF64VuCPmwAWAGqp4CNlDBZk1YkgMBmZtgAcSAJyvCLnw-GkoCGsLFRxUBl4cAwQAehk5YAAPKETEAxcoB2woQhzTRAsMMBwLNFxEADEALigAeXZtfBAAewAbVtwQNGboBQgwXl1YHqgAQygkbKNA3FbRgE8oFrBWqBG5oXZY+Plk1IR0qAA2Axy82ALUIspo9HKEABUAC1xM6gcoaYBxAAUASVqUAwAHcwGgQC8MGMVFAAHI-WHoCCjFRFDDdaAACgYuDQ4zgE0IhAAlOgnhBmgBXADmTygACUAPoAIWoUGer1azVGCiBs06GKWT2azQwuChEFmCyWzRWa2gG3Y1G2sl2KTSGR8ZwsZkKxWQtzKiAAwlyxVAAOQwVrzBijEAAawt6GaYwUvLQL1gkoIwJeszAKmpUAUYGaClwAG41g0YRaGPN2KMMBhnd0oOxKWgxlBDJ8oFodNBA3z+uKoKDPZY0Mq4qrgPT1QciRNXFADC33Ed8uZUMDdcJRolIsg6HWEsCmx3cpMMgYTq2XD22OP5JP9vPTrPjnnZ8uVQkAIJCIS4oYblvb6cTYjdi4IBzCe6FCi91fAVYXme39stww6sIoFaApqBQd9oC-XMO1-XIHEIW4B0CNIiQPNcp0vH8HAcI4bzvAd33XRAsJwuDTmvUj93fEAm2I3DMmwzJ4J7OgljEIh3FQpIm2wbwoNyfj9yBfV7FQI0EElOZFmWVpOL2EoMmwBcBPvPUblE59QKgAAJAAaGC6ieIROMI5tYPg9sFwolTON5DdLPMzc+MEoRmkoWEDHsJd3w0KcFyQH9HP-ZzXIQQtdH0YxjN805-ygZwFI7SidmAEz4roxz-Lwy53yhfY0titLMsE9hLigFzsuSuTuAMGKkBw5TEOuA11MQBMkxTTjcrSPyAr81xisoYQQpXSqpxcGcjgyBrSqakS7kQa1bXtB1OIgxAjl3NsNtOKzEIEBBsA498qpOXJsDq7VBNmw17kVVt8xiOKoAAdU4gBlJtmEmfSvuyYiBt7MrhtQVp9VKJ4wgIqcu3onCDBhrIALfZKTPhkieP0xGkvrVGW0RzGmOsqG8rbWi0qxomUc+0mMe2zIeOxicmzp-GWcJ-DkuzDdDhqjJfvbdJitCZBwnuKgLHKhAIxgUZKVaNAory77HGyfmPKy5Gcc+5WLrVlDKfrVoYSV7J1bigK6p7G6TU44MTYs041a7Rn5EbfZsGYOHDh4udBesiszBYmB9p4iWhxHJBqCOZhwjAZNKCjmPwO4z2Bd99PBP28WRvrHkACsoHuBwl0LsHUG0RB6XFKw8RUEBXkxZpDIwYlhGF8XhH2bOiWoQo0EQZlk1eSU0UFJlmRB5pqUG9988LoiAlwMvC+FipJQAR0pbQQEWTEtI0VuyvbiWu6JPvEDeysIUDYNg9wTft-mSfp5F99XQrh8Y6BEqEBOZiOEuDsBI7AoAjhfCwHOCQYjNGSu-YWd4MA-z-veFAGZBpAPkCAkcj4gSvhoIQJOyVoGwJXkRFiiD1oIVKj-OgGDgBYMQAEDAr4wJEJgdIYAMRwAQBkEAA}
\def\linkG{https://www.falstad.com/circuit/circuitjs.html?ctz=DwYwlgTgBAZgvAIgIwKgFwM6KQUwLRIpRgjYAMAdAEyogBGiArGagA4MLOoBuETqAW0zZ8hBAD4AUFCjABAqADsBiAOaswAezzLNGPDACGGNKngIWUbmkSUAbKjo5bFMgGZUGDpYAmGAJ62jj7YjjYWnnQhEQD0UjJyCspqGtqsAnoGxqawQVbh9o7OFq5UjJF5foERUFGhtQWR0Sxx0rIAHlDtiAAsABxQblRQZCNmeRhgiFSWaMUA4gBcUABi2VCsAK4YABZ4WwA2B1A4Pqo4GFCailCGUMxQAPI7CgAKAMIAQugQhiA4CFaCU63QQ-Sgdh6IzGuRqk2wPVQc0QABUdjgfn8MQAKJBKS4+HAHQz+AA0tzomk2aCgSEYUBACmuqwASj0AJQnRR+W40tDo4iKLY0zQwW6KTT8nDQEA7MCsZb01grDYHbZQACiAGUACJQHxaQkYCiA+IdLq9AZ9AajSzmSzw5CI9DFR7U4WC4iYWDZC40jA4NBoMCKVRQbEsgD6zMMlOptKoBHpzwE7IoUC1BzAhIgGGWAYgYAuUAgFzAJk00CjWvJGAOOAA7iXDHN01qQJpWDhlq8AJIAOXFPigLIAGiagbIWRaEABOOxQMoL8FL8Y1BvmGiwPjIMh7-eUSwCQztbjYCh2cq1MDGc+X01tYBN0HzxeMBdIKgDVewlpm4C9jOn7fu+tIAOwLj+9qeAc-BQHKiBuBQHiTsAYqglQPSjDMAxuGQUI4WuRAZCkWj7Bk+hGCYD4JOh0xYYM+GMVCeHOtBUAkQg6hkbolHZDRsjPvR2FkN+DFQUEqFCQgmHYaBsnMWuf6PtJrGLphzFvg4v4CU+M5qUumlfn0Sm6WGoJqUMoyWYwW7Qah04WUxSCzsMalIHYdp5BueTHqeiHIbp0lIBBWmLuJ76mVJ+lMYZbhuJBkU6ahBw3BZCVhfFkGYWuzrIggfb9rp0DpYly5DFppkli4KCoWgxygj0FWGXY1qVTp+rFLgeBkLOqBgConCWA2UwIEMZkzk1wzGVAU2DDMUX-uZiCtd+bWrfNXmxP+0kbYZYG2e19k7TOB3TaBfQeUdkknaCl1lVA1pydpx2Pqlp2Hftn05bCeXFGOukgDOT1hSDWVVaQnCiDQqHLQgYMZWDrWLY+w53Vdhn3bSrlVZxhJGJsBymA5wNXS5wxYx5W3DexflnmNgWoWjiBY1Zj0Y0ltOaIg+OGITxNLaTC5s6ztko8Cp0ZVhtowuxjpIH9iClv8YDcNKgrCrphg+AAVic2BgagOCjZYOCKIgLJliYhiKP84aaC8GDsoI5slMRoJIHunjhFq0pFpcpaTBWVaRlqWu6-rnBGybRuu5mjYji2OIAGoxIoTsuzVgigleJiIG6aAenWCe-HM4ap+nztQAcmiqKEqGaFy2DaRgrArWuW5txEZoJKwVgTJD5STr3kc-fTDqD7pMSaP+jdm9MhtQK37ewkQXcoD3sh9-THhL5Dh5D5vwB98UVB9f3Y2ePvriH4+09SMAMTgBAUhAA}


% numbered tasks: P = preparation (paper), S = simulation, D = discussion
\newcounter{task}[section]
\renewcommand{\thetask}{\thesection.\arabic{task}}
\newcommand{\task}[1]{\par\medskip\refstepcounter{task}%
  {\sffamily\bfseries\color{rwuviolet}Task \thetask}\ {\sffamily\small\color{rwucyandark}[#1]}\quad}
\newtcolorbox{modelbox}[1]{colback=rwucyan40!25,colframe=rwucyan,fonttitle=\sffamily\bfseries,title=#1,boxrule=0.6pt,arc=1mm}
\newtcolorbox{solbox}{colback=rwuviolet!6,colframe=rwuviolet,fonttitle=\sffamily\bfseries,title=Solution,boxrule=0.6pt,arc=1mm,breakable}
\ifdefined\solutions
  \newenvironment{solution}{\begin{solbox}}{\end{solbox}}
\else
  \usepackage{comment}
  \excludecomment{solution}
\fi
% value measured in the simulation: empty cell for students, filled in the solutions
\ifdefined\solutions
  \newcommand{\simv}[1]{\textcolor{rwuviolet}{#1}}
\else
  \newcommand{\simv}[1]{}
\fi
% table of given (calculated) values; simulation columns are filled in by the students
\newenvironment{results}[2]{\par\smallskip\begin{center}\small\sffamily
  {\color{rwucyandark}\bfseries #1}\\[1mm]\renewcommand{\arraystretch}{1.25}\begin{tabular}{#2}\toprule}
  {\bottomrule\end{tabular}\end{center}}
% "open circuit X" link with the file name
\newcommand{\circuitlink}[2]{\href{#1}{\color{rwucyandark}\sffamily\bfseries $\blacktriangleright$\ open in Falstad}\ \ (\texttt{#2})}
\newcommand{\schematic}[2][0.62]{\begin{center}\includegraphics[width=#1\textwidth]{figures/falstad/#2_schematic}\end{center}}

\begin{document}

\begin{center}
  \rwulogo\\[6mm]
  {\sffamily\bfseries\LARGE\color{rwuviolet} GPIO Output Stages in Simulation}\\[2mm]
  {\sffamily\Large\color{rwucyandark} Push-pull, open drain and open collector: switching speed, voltage and current limits}\\[3mm]
  {\sffamily Embedded Computing -- Master of Engineering -- Worksheet\ifdefined\solutions\ with solutions\fi}
\end{center}

\section*{Overview}

Every GPIO pin ends in an output stage: two transistors (push-pull) or one transistor to ground (open drain,
open collector). In this worksheet you will take these stages apart in a circuit simulator. You will find out
how much current a pin can deliver, how fast it can switch a load, why a pull-up resistor makes a bus slow,
and which voltages and currents destroy a pin or the circuit behind it.

\textbf{Learning outcomes.} After this worksheet you can
\begin{itemize}
  \item model a GPIO output stage by on-resistance and saturation current and relate the model to the
        $V_{OH}$/$V_{OL}$ specifications of a data sheet,
  \item predict rise and fall times, the maximum toggle frequency and the dynamic power of a push-pull output,
  \item dimension the pull-up resistor of an open-drain bus such as I$^2$C from rise-time and $V_{OL}$ limits,
  \item judge where the maximum output voltage and current of a pin come from (body diodes, injection current,
        back-powering, transistor breakdown and saturation) and protect a driver against inductive kick-back,
  \item recognise when a PCB trace has to be treated as a transmission line and terminate it.
\end{itemize}

\textbf{Organisation.} About 4 hours, in teams of two. You do not have to calculate: all results of the
model are given in tables. Tasks marked \textsf{[P]} (preparation) ask you to read these values and understand
where they come from \emph{before} you open the simulation. \textsf{[S]} means simulate, measure and enter your
values in the empty column \emph{Simulation} of the tables. \textsf{[D]} means discuss and explain. Hand in a short report (PDF): for each task the completed tables, screenshots of the
scopes and your explanation. When a measurement differs from the model value by
more than about \SI{10}{\percent}, explain why.

\textbf{Tool.} Falstad CircuitJS runs in the browser: \url{https://www.falstad.com/circuit/}. Each part of this
sheet has a link that opens the prepared circuit. The circuit files are also in the folder
\href{https://github.com/mopfeil/EmbeddedComputing/tree/main/examples/interfaces/falstad}{\texttt{examples/interfaces/falstad/}}
of the lecture repository (open them with \emph{File $\rightarrow$ Open File}).

\begin{modelbox}{Working with CircuitJS}
\begin{itemize}[leftmargin=*]
  \item Hover over an element to see its voltages, currents and power in the lower right corner.
  \item Click on a logic input (\textsf{H}/\textsf{L}) or a switch to toggle it. Double-click an element to edit
        its value. Sliders on the right change the values marked in the circuit.
  \item Scopes: right-click $\rightarrow$ \emph{Properties} sets the time base (``scope speed''), the scale and
        shown values (e.g.\ average, peak-to-peak). Drag with the mouse in a scope to measure time differences.
  \item \emph{Reset} restarts the simulation, \emph{Run/STOP} freezes it for reading values.
  \item \emph{Options $\rightarrow$ Other Options} lets you change the simulation time step. If edges look
        jagged, make the time step smaller.
\end{itemize}
\end{modelbox}

\begin{modelbox}{The transistor model used in all circuits}
The MOSFETs follow the square-law (Shichman--Hodges) model with threshold $V_T=\SI{0.6}{V}$ and
\[
  I_D=\begin{cases}
     \beta\left[(V_{GS}-V_T)\,V_{DS}-\tfrac12 V_{DS}^2\right] & V_{DS} < V_{GS}-V_T \quad\text{(triode)}\\[1mm]
     \tfrac{\beta}{2}\,(V_{GS}-V_T)^2 & V_{DS} \ge V_{GS}-V_T \quad\text{(saturation)}
  \end{cases}
\]
with $\beta_N=\SI{12}{mA/V^2}$ for the NMOS and $\beta_P=\SI{8}{mA/V^2}$ for the PMOS of the output stage (PMOS:
use $V_{SG}$, $V_{SD}$). Both have a body diode. The supply is $V_{DD}=\SI{3.3}{V}$.
The model knows no gate or junction capacitances, no temperature and no breakdown. Pin and load capacitances are
drawn as explicit capacitors. Keep in mind what the model cannot show when you interpret the results.
\end{modelbox}

%======================================================================
\needspace{0.45\textheight}
\section{Push-pull output: static characteristic}
%======================================================================
\circuitlink{\linkA}{A\_pushpull\_characteristic.txt}

The logic input is the bit in the output register. Two inverters (the predriver and the CMOS stage) make the pin
follow it. A triangle voltage source acts as a curve tracer: it forces the pin voltage from \SI{-0.5}{V} to
\SI{3.8}{V}, and an ammeter measures the pin current (positive = current out of the pin). The left scope shows the
current over the voltage, the right scope the power in both transistors.

\schematic[0.55]{A}

\task{P} Table~1.1 gives the on-resistance (valid for small $V_{DS}$) and the saturation current of both
transistors for $V_{GS}=V_{DD}$. Which one is the stronger transistor, and why do chip designers make the PMOS
larger than the NMOS?

\begin{results}{Table 1.1: model values of the output transistors ($V_{GS}=\SI{3.3}{V}$)}{lccc}
  & formula & NMOS & PMOS\\ \midrule
  $\beta$ & & \SI{12}{mA/V^2} & \SI{8}{mA/V^2}\\
  $R_{on}$ & $1/(\beta(V_{GS}-V_T))$ & \SI{30.9}{\ohm} & \SI{46.3}{\ohm}\\
  $I_{sat}$ & $\tfrac{\beta}{2}(V_{GS}-V_T)^2$ & \SI{43.7}{mA} & \SI{29.2}{mA}\\
\end{results}

\begin{solution}
$R_{on,N}=1/(\SI{12}{mA/V^2}\cdot\SI{2.7}{V})=\SI{30.9}{\ohm}$, $R_{on,P}=1/(\SI{8}{mA/V^2}\cdot\SI{2.7}{V})=\SI{46.3}{\ohm}$.
$I_{sat,N}=\tfrac12\cdot 12\cdot 2.7^2\,\si{mA}=\SI{43.7}{mA}$, $I_{sat,P}=\SI{29.2}{mA}$.
The NMOS is stronger. Hole mobility is about 2--3 times lower than electron mobility, so a PMOS needs a
correspondingly wider channel for the same current. Here the PMOS is still weaker, which is typical: data sheets
often allow a larger $V_{OH}$ drop than $V_{OL}$ rise, or specify asymmetric currents.
\end{solution}

\task{S} Set the output to \textsf{H} and then to \textsf{L}. From the curve, read $V_{OH}$ and $V_{OL}$ at the
currents of Table~1.2 (Run/STOP and hovering over the PIN label helps) and enter them next to the model values.
Also read the short-circuit currents (pin at \SI{0}{V} for \textsf{H}, pin at \SI{3.3}{V} for \textsf{L}).

\begin{results}{Table 1.2: output levels over the load current}{lcccccc}
  $|I|$ & \SI{2}{mA} & \SI{4}{mA} & \SI{8}{mA} & \SI{12}{mA} & \SI{20}{mA} & short circuit\\ \midrule
  $V_{OH}$, model & \SI{3.21}{V} & \SI{3.11}{V} & \SI{2.90}{V} & \SI{2.67}{V} & \SI{2.11}{V} & \SI{29.2}{mA}\\
  $V_{OH}$, simulation & \simv{3.21} & \simv{3.11} & \simv{2.90} & \simv{2.67} & \simv{2.11} & \simv{29.2}\\ \midrule
  $V_{OL}$, model & \SI{0.06}{V} & \SI{0.13}{V} & \SI{0.26}{V} & \SI{0.40}{V} & \SI{0.71}{V} & \SI{43.7}{mA}\\
  $V_{OL}$, simulation & \simv{0.06} & \simv{0.13} & \simv{0.26} & \simv{0.40} & \simv{0.71} & \simv{43.7}\\
\end{results}

\begin{solution}
Model and simulation agree to within the reading accuracy:\\[1mm]
\begin{tabular}{lccccc}
  \toprule
  $|I|$ / mA & 2 & 4 & 8 & 12 & 20\\
  \midrule
  $V_{OH}$ / V & 3.21 & 3.11 & 2.90 & 2.67 & 2.11\\
  $V_{OL}$ / V & 0.06 & 0.13 & 0.26 & 0.40 & 0.71\\
  \bottomrule
\end{tabular}\\[1mm]
Short-circuit currents: \SI{29.2}{mA} (high, PMOS saturated), \SI{43.7}{mA} (low, NMOS saturated).
\end{solution}

\task{D} The RP2040 data sheet specifies $V_{OH}\ge\SI{2.62}{V}$ and $V_{OL}\le\SI{0.5}{V}$ at a load current
equal to the selected drive strength (2, 4, 8 or \SI{12}{mA}) and $IOVDD=\SI{3.3}{V}$. The model delivers
\SI{12.8}{mA} at $V_{OH}=\SI{2.62}{V}$ and sinks \SI{12}{mA} at $V_{OL}=\SI{0.40}{V}$. Which drive strength setting
does the model represent? Drive strength is implemented by switching transistor ``fingers'' in parallel, i.e.\ by
scaling $\beta$ ($I_D\propto\beta$): the \SI{4}{mA} setting needs $\beta_P=\SI{2.5}{mA/V^2}$. What fraction of
the fingers is that, and what does it mean for the switching speed (part 3)?

\begin{solution}
At $V_{OH}=\SI{2.62}{V}$ the model PMOS delivers \SI{12.8}{mA}, the NMOS sinks \SI{12}{mA} at $V_{OL}=\SI{0.4}{V}$:
the model corresponds to the \SI{12}{mA} setting. Since $I_D\propto\beta$ at a fixed operating point,
$\beta_P=\SI{8}{mA/V^2}\cdot 4/12.8=\SI{2.5}{mA/V^2}$ meets the \SI{4}{mA} setting, i.e.\ about one third of
the fingers.
\end{solution}

\task{S/D} A red LED ($V_F\approx\SI{2.0}{V}$) with a \SI{100}{\ohm} series resistor is connected from the pin to
GND. The naive calculation gives $I=(\SI{3.3}{V}-\SI{2.0}{V})/\SI{100}{\ohm}=\SI{13}{mA}$. Draw the load line
$I=(V_{PIN}-V_F)/R$ into your measured characteristic (it passes through \SI{2.0}{V}/\SI{0}{mA} and
\SI{3.3}{V}/\SI{13}{mA}) and read the operating point. The model gives $V_{PIN}=\SI{2.86}{V}$ and $I=\SI{8.6}{mA}$.
Why is the current a third lower than the naive value? Optional: verify it by replacing the sweep source with an
LED and a resistor (\emph{Draw $\rightarrow$ Outputs and Labels $\rightarrow$ Add LED}).

\begin{solution}
Intersection with the PMOS triode curve: $V_{SD}=\SI{0.44}{V}$, i.e.\ $V_{PIN}=\SI{2.86}{V}$ and
$I=\SI{8.6}{mA}$, a third less than the naive value. The output resistance of the pin is not negligible
compared to \SI{100}{\ohm}; for a defined LED current either the resistor must be larger (pin resistance becomes
negligible) or the output stage must be included in the calculation.
\end{solution}

\task{D} Look at the curve below \SI{0}{V} and above \SI{3.3}{V}. Which element conducts there? Relate this to the
absolute maximum rating ``pin voltage $-0.5\ldots IOVDD+\SI{0.5}{V}$'' of the RP2040 and to the term
\emph{injection current}. Why can a large injection current upset an analogue measurement or even a digital
circuit on the same chip?

\begin{solution}
Above $V_{DD}$ the body diode of the PMOS (and in a real chip also the ESD clamp diode) conducts into the supply;
below \SI{0}{V} the body diode of the NMOS conducts from GND. From about \SI{0.4}{V} beyond the rails the current
rises exponentially (in the model to about \SI{9}{mA} at \SI{3.7}{V}). The limits $\pm\SI{0.5}{V}$ keep these
diodes nearly off. Injected current flows as minority carriers into the substrate or into the supply rail: it can
shift the ADC reference and neighbouring channels, trigger the parasitic thyristor of the CMOS structure
(latch-up), or raise $V_{DD}$ (see part~\ref{sec:overvoltage}).
\end{solution}

\task{S/D} Read the maximum power in the transistors from the right scope (model: \SI{96}{mW} in the PMOS for a
short circuit to GND, \SI{110}{mW} at \SI{-0.5}{V}). Many microcontrollers allow \SIrange{50}{100}{mA} in total for
all pins of one supply domain. What does this mean for the number of pins that may drive LEDs at the same time,
and what happens to the chip temperature if a pin is shorted to GND?

\begin{solution}
In a short circuit to GND the PMOS dissipates $\SI{3.3}{V}\cdot\SI{29}{mA}\approx\SI{96}{mW}$ (up to
\SI{110}{mW} at $-0.5$\,V). One pin is survivable, but the total current budget (e.g.\ \SI{50}{mA}) allows only
about five LEDs at \SI{8.6}{mA}. The power is dissipated in a tiny area at the pad; together with metal
migration in the supply lines this is why per-pin and per-domain current limits exist. Larger loads need an
external driver (part~\ref{sec:oc}).
\end{solution}

%======================================================================
\needspace{0.45\textheight}
\section{Two push-pull outputs on one line}
%======================================================================
\circuitlink{\linkC}{C\_bus\_contention.txt}

\schematic[0.62]{C}

\task{P/S} Output A drives high, output B drives low. The line voltage follows from $I_{D,P,A}(V)=I_{D,N,B}(V)$;
both transistors work in the triode region. Measure the values of Table~2.1 in the simulation.

\begin{results}{Table 2.1: contention between two push-pull outputs}{lcccc}
  & line voltage & current & power PMOS A & power NMOS B\\ \midrule
  model & \SI{1.09}{V} & \SI{28.2}{mA} & \SI{62}{mW} & \SI{31}{mW}\\
  simulation & \simv{1.09} & \simv{28.2} & \simv{62.3} & \simv{30.8}\\
\end{results}

\begin{solution}
Both transistors are in the triode region. $\beta_P[2.7\,(3.3-V)-\tfrac12(3.3-V)^2]=\beta_N[2.7\,V-\tfrac12V^2]$
gives $V=\SI{1.09}{V}$ and $I=\SI{28.2}{mA}$. Power: $P_{P,A}=(3.3-1.09)\,\si{V}\cdot\SI{28.2}{mA}=\SI{62}{mW}$,
$P_{N,B}=\SI{31}{mW}$.
\end{solution}

\task{D} What logic level does a receiver see on the line?

\begin{solution}
Simulation: \SI{28.20}{mA}, \SI{1.09}{V}, \SI{62.3}{mW} and \SI{30.8}{mW}. \SI{1.09}{V} lies between
$V_{IL,max}$ and $V_{IH,min}$ of a \SI{3.3}{V} CMOS input (about 0.8--2.0\,V): the level is undefined.
\end{solution}

\task{D} Name three situations in embedded practice where this happens. How do tri-state outputs, open-drain
outputs and series resistors prevent the damage? (A series resistor of about \SI{290}{\ohm} in each output
limits the contention current to \SI{5}{mA}: $2R+R_{on,P}+R_{on,N}=\SI{3.3}{V}/\SI{5}{mA}$. What does it cost?)

\begin{solution}
Examples: two masters on a shared line (bus arbitration error), a bidirectional data line where both sides
drive at the same time during the turn-around, bit-banged I$^2$C with the pin configured push-pull, a debugger
and the board both driving reset, two boards connected with swapped TX lines, a pin shorted to its neighbour.
Tri-state: only one driver is enabled at a time (protocol must guarantee it). Open drain: nobody drives high, so a
conflict only produces a low level. Series resistors: $2R+R_{on,P}+R_{on,N}\ge \SI{3.3}{V}/\SI{5}{mA}=\SI{660}{\ohm}$
$\Rightarrow R\approx\SI{290}{\ohm}$ per output (\SI{330}{\ohm} standard value), at the cost of speed (part 3).
\end{solution}

%======================================================================
\needspace{0.45\textheight}
\section{Push-pull output: switching speed and dynamic power}
%======================================================================
\circuitlink{\linkB}{B\_pushpull\_capacitive\_load.txt}

The square wave is the program toggling the pin. The predriver is a logic inverter with an adjustable slew rate
(how fast the gate voltage changes). The output drives the load capacitance $C_L$ (pin, trace, inputs of other
chips). An ammeter measures the current taken from the supply.

\schematic[0.55]{B}

\task{P} Table~3.1 gives the rise and fall times (\SIrange{10}{90}{\percent}) from two models: (a) the RC
approximation $t_r=2.2\,R_{on}C_L$, (b) the square-law model, in which the transistor starts in saturation
(constant current $I_{sat}$) and enters the triode region when $V_{DS}=V_{GS}-V_T$ (integration of
$dt=C_L\,dV/I_D(V)$). Why is the RC approximation too optimistic?

\begin{solution}
(a) $t_r=2.2\cdot\SI{46.3}{\ohm}\cdot C_L$, $t_f=2.2\cdot\SI{30.9}{\ohm}\cdot C_L$: \SI{1.0}{ns}/\SI{0.68}{ns},
\SI{10.2}{ns}/\SI{6.8}{ns}, \SI{102}{ns}/\SI{68}{ns}.\\
(b) Rising edge, $a=V_{GS}-V_T=\SI{2.7}{V}$: from \SI{0.33}{V} to \SI{0.6}{V} the PMOS is saturated,
$t_1=C_L\cdot\SI{0.27}{V}/\SI{29.2}{mA}$; then triode from $V_{SD}=\SI{2.7}{V}$ down to \SI{0.33}{V}:
\[
  t_2=\frac{C_L}{\beta}\int_{0.33}^{2.7}\frac{dx}{x(a-x/2)}=\frac{C_L}{\beta a}\left[\ln\frac{x}{2a-x}\right]_{0.33}^{2.7}
     =\frac{C_L}{\beta}\cdot\SI{1.012}{V^{-1}}.
\]
Result $t_r=\SI{1.36}{ns}$, \SI{13.6}{ns}, \SI{136}{ns} and $t_f=\SI{0.90}{ns}$, \SI{9.0}{ns}, \SI{90}{ns}.
The RC estimate is about \SI{25}{\percent} too optimistic, because at the start of the edge the transistor
delivers only $I_{sat}$ instead of $V_{DD}/R_{on}$.
\end{solution}

\task{S} Set the predriver slew rate to its maximum and measure $t_r$ and $t_f$ for the three capacitances (set the
scope time base accordingly, or zoom in with Run/STOP). For which $C_L$ does the predriver dominate the edge?

\begin{results}{Table 3.1: rise and fall times at the pin}{lcccccc}
  & \multicolumn{2}{c}{$C_L=\SI{10}{pF}$} & \multicolumn{2}{c}{$C_L=\SI{100}{pF}$} & \multicolumn{2}{c}{$C_L=\SI{1}{nF}$}\\
  & $t_r$ & $t_f$ & $t_r$ & $t_f$ & $t_r$ & $t_f$\\ \midrule
  (a) RC model & \SI{1.0}{ns} & \SI{0.68}{ns} & \SI{10.2}{ns} & \SI{6.8}{ns} & \SI{102}{ns} & \SI{68}{ns}\\
  (b) square-law model & \SI{1.36}{ns} & \SI{0.90}{ns} & \SI{13.6}{ns} & \SI{9.0}{ns} & \SI{136}{ns} & \SI{90}{ns}\\
  simulation & \simv{1.5} & \simv{1.1} & \simv{13.6} & \simv{9.1} & \simv{136} & \simv{90}\\
\end{results}

\begin{solution}
Simulation (slew \SI{2}{V/ns} or more): \SI{1.5}{ns}/\SI{1.1}{ns}, \SI{13.6}{ns}/\SI{9.1}{ns},
\SI{136}{ns}/\SI{90}{ns}, i.e.\ exactly model (b) for \SI{100}{pF} and \SI{1}{nF}. At \SI{10}{pF} and below
the gate itself needs about \SI{1.5}{ns} to swing, so the predriver dominates.
\end{solution}

\task{S/D} Increase the toggle frequency with $C_L=\SI{1}{nF}$ until the pin no longer reaches valid logic levels
($V_{IH}=0.7\,V_{DD}$, $V_{IL}=0.3\,V_{DD}$). Compare with the rule of thumb $f_{max}\approx 1/(2\,t_r)\approx
\SI{3.7}{MHz}$ (slower edge). Why do data sheets specify the maximum GPIO frequency only together with a load
capacitance?

\begin{solution}
Each half period must be long enough for the output to cross the threshold: roughly $f_{max}\approx
1/(2\,t_r)$ with the slower edge, here $\approx\SI{3.7}{MHz}$ for \SI{1}{nF}; beyond that the swing shrinks and
at about \SIrange{5}{6}{MHz} the levels are no longer valid. Since $t_r\propto C_L$, a frequency without a
load capacitance is meaningless (typical: ``\SI{50}{MHz} at \SI{30}{pF}'').
\end{solution}

\task{S} Set $f=\SI{1}{MHz}$ and the fastest slew rate. Read the average supply current (right scope) for the three
capacitances and compare with $I=C_LV_{DD}f$. The energy taken from the supply per cycle is $C_LV_{DD}^2$,
independent of $R_{on}$. Where is it dissipated?

\begin{results}{Table 3.2: average supply current at \SI{1}{MHz}}{lccc}
  & $C_L=\SI{10}{pF}$ & $C_L=\SI{100}{pF}$ & $C_L=\SI{1}{nF}$\\ \midrule
  $C_LV_{DD}f$ & \SI{33}{\micro A} & \SI{330}{\micro A} & \SI{3.3}{mA}\\
  simulation & \simv{34} & \simv{330} & \simv{3300}\\
\end{results}

\begin{solution}
Simulation: \SI{34}{\micro A}, \SI{330}{\micro A}, \SI{3300}{\micro A}; formula: \SI{33}{\micro A},
\SI{330}{\micro A}, \SI{3300}{\micro A}. During charging the supply delivers $Q V_{DD}=C_LV_{DD}^2$; half of it is
stored in $C_L$, the other half is converted to heat in the PMOS, whatever its resistance. During discharge the
stored half is dissipated in the NMOS. $P=C_LV_{DD}^2f$ is the dynamic power of all CMOS logic.
\end{solution}

\task{S/D} Set $C_L=\SI{1}{pF}$ and reduce the predriver slew rate to \SI{0.02}{V/ns}. Observe the supply current.
Both transistors conduct while $V_T<V_G<V_{DD}-V_T$ (\SIrange{0.6}{2.7}{V}), at \SI{0.02}{V/ns} that is
\SI{105}{ns} per edge. Measure the average supply current and compare it with $C_LV_{DD}f=\SI{3.3}{\micro A}$
(simulation: \underline{\hspace{2cm}}\simv{376\,µA}). Why do microcontroller input pins have Schmitt triggers, and
why does a slowly changing analogue voltage on a digital input increase the current consumption?

\begin{solution}
Both transistors conduct while $V_T<V_G<V_{DD}-V_T$, i.e.\ \SIrange{0.6}{2.7}{V}; at \SI{0.02}{V/ns} this takes
\SI{105}{ns} per edge. A cross-conduction (shoot-through) current flows directly from $V_{DD}$ to GND (peak
\SI{5.8}{mA}). The average supply current rises from \SI{6}{\micro A} to \SI{376}{\micro A}, more than
100 times $C_LV_{DD}f=\SI{3.3}{\micro A}$. An input stage is also a CMOS inverter: a slow input edge or an analogue
voltage near $V_{DD}/2$ keeps both transistors on (static cross current) and may make the output oscillate.
A Schmitt trigger switches quickly and only once; unused or analogue pins should have the digital input buffer
disabled.
\end{solution}

%======================================================================
\needspace{0.45\textheight}
\section{Open-drain bus versus push-pull}
%======================================================================
\circuitlink{\linkD}{D\_open\_drain\_bus.txt}

Two open-drain devices share a line with one pull-up resistor $R_P$, as on an I$^2$C bus. Device~1 sends a clock;
device~2 can hold the line low. A push-pull stage drives the same capacitance for comparison.

\schematic[0.62]{D}

\task{P} The I$^2$C rise time is measured from \SI{30}{\percent} to \SI{70}{\percent} of $V_{DD}$. For the RC charge
$V(t)=V_{DD}(1-e^{-t/\tau})$ it is $t_r=\tau\ln(0.7/0.3)=0.8473\,R_PC_b$, i.e.\ \SI{796}{ns} for
$R_P=\SI{4.7}{k\ohm}$ and $C_b=\SI{200}{pF}$.

\begin{solution}
$V(t)=V_{DD}(1-e^{-t/\tau})$ reaches $0.3V_{DD}$ at $t_{30}=\tau\ln(1/0.7)$ and $0.7V_{DD}$ at $t_{70}=\tau\ln(1/0.3)$;
the difference is $\tau\ln(0.7/0.3)=0.8473\,\tau$. With \SI{4.7}{k\ohm} and \SI{200}{pF}: \SI{796}{ns}.
\end{solution}

\task{S} Measure the rise and fall times of BUS and the rise time of PP (same thresholds). Explain the difference
between rising and falling edge of the open-drain line.

\begin{results}{Table 4.1: edges with $R_P=\SI{4.7}{k\ohm}$, $C_b=\SI{200}{pF}$ (\SIrange{30}{70}{\percent})}{lccc}
  & $t_r$ BUS & $t_f$ BUS & $t_r$ PP\\ \midrule
  model (RC) & $0.8473\,R_PC_b=\SI{796}{ns}$ & $0.8473\,R_{on,N}C_b\approx\SI{5}{ns}$ & $0.8473\,R_{on,P}C_b\approx\SI{8}{ns}$\\
  simulation & \simv{796} & \simv{8} & \simv{12}\\
\end{results}

\begin{solution}
BUS: $t_r=\SI{796}{ns}$ (exactly the prediction), $t_f=\SI{8}{ns}$; PP: $t_r=\SI{12}{ns}$. The falling edge is
driven by the NMOS (\SI{31}{\ohm}), the rising edge only by the pull-up resistor (\SI{4.7}{k\ohm}): the ratio of
the resistances is the ratio of the edges. The asymmetry is the price for the wired-AND function.
\end{solution}

\task{P} Table~4.2 gives the permitted range of $R_P$ for $C_b=\SI{200}{pF}$ and $V_{DD}=\SI{3.3}{V}$ according to
the I$^2$C specification (NXP UM10204): $V_{OL}\le\SI{0.4}{V}$ at a sink current of \SI{3}{mA}, $t_r\le\SI{1000}{ns}$
(Standard-mode, \SI{100}{kHz}) and $t_r\le\SI{300}{ns}$ (Fast-mode, \SI{400}{kHz}). Which limit comes from which
requirement?

\begin{results}{Table 4.2: pull-up range for $C_b=\SI{200}{pF}$}{lccc}
  & $R_{P,min}=\dfrac{V_{DD}-V_{OL}}{I_{OL}}$ & $R_{P,max}=\dfrac{t_r}{0.8473\,C_b}$ & standard values\\ \midrule
  Standard-mode & \SI{967}{\ohm} & \SI{5.9}{k\ohm} & \SI{1}{k\ohm}\ldots\SI{5.6}{k\ohm}\\
  Fast-mode & \SI{967}{\ohm} & \SI{1.77}{k\ohm} & \SI{1}{k\ohm}\ldots\SI{1.5}{k\ohm}\\
\end{results}

\begin{solution}
$R_{P,min}=(V_{DD}-V_{OL})/I_{OL}=\SI{2.9}{V}/\SI{3}{mA}=\SI{967}{\ohm}$.
$R_{P,max}=t_r/(0.8473\,C_b)$: \SI{5.9}{k\ohm} (Standard-mode), \SI{1.77}{k\ohm} (Fast-mode).
Standard-mode: \SI{1}{k\ohm}\ldots\SI{5.6}{k\ohm}; Fast-mode: \SI{1}{k\ohm}\ldots\SI{1.8}{k\ohm} (standard values,
\SI{1.8}{k\ohm} is just at the limit).
\end{solution}

\task{S} Check the limits in the simulation: complete Table~4.3. Then increase $C_b$ to \SI{400}{pF} (the maximum
allowed by I$^2$C) with $R_P=\SI{4.7}{k\ohm}$ (model: $t_r=\SI{1.59}{\micro s}$): which mode still works?

\begin{results}{Table 4.3: low level and rise time over $R_P$ ($C_b=\SI{200}{pF}$)}{lcccccc}
  $R_P$ & \SI{5.9}{k\ohm} & \SI{4.7}{k\ohm} & \SI{1.8}{k\ohm} & \SI{1}{k\ohm} & \SI{470}{\ohm} & \SI{220}{\ohm}\\ \midrule
  $I_{low}$, model & \SI{0.56}{mA} & \SI{0.70}{mA} & \SI{1.8}{mA} & \SI{3.2}{mA} & \SI{6.6}{mA} & \SI{13.1}{mA}\\
  $V_{OL}$, model & \SI{17}{mV} & \SI{22}{mV} & \SI{56}{mV} & \SI{99}{mV} & \SI{204}{mV} & \SI{406}{mV}\\
  $V_{OL}$, simulation & \simv{17} & \simv{22} & \simv{56} & \simv{101} & \simv{211} & \simv{437}\\
  $t_r$, model & \SI{1000}{ns} & \SI{796}{ns} & \SI{305}{ns} & \SI{169}{ns} & \SI{80}{ns} & \SI{37}{ns}\\
  $t_r$, simulation & \simv{1000} & \simv{796} & \simv{306} & \simv{168} & \simv{78} & \simv{36}\\
\end{results}

\begin{solution}
$R_P=\SI{1.8}{k\ohm}$: $t_r=\SI{306}{ns}$, just above the Fast-mode limit (\SI{1.5}{k\ohm} gives margin);
$R_P=\SI{5.9}{k\ohm}$: $t_r=\SI{1000}{ns}$.\\[1mm]
\begin{tabular}{lcccc}
  \toprule
  $R_P$ & \SI{4.7}{k\ohm} & \SI{1}{k\ohm} & \SI{470}{\ohm} & \SI{220}{\ohm}\\
  \midrule
  $I_{low}$ / mA & 0.70 & 3.2 & 6.6 & 13.0\\
  $V_{OL}$ / mV & 22 & 101 & 211 & 437\\
  $t_r$ / ns & 796 & 168 & 78 & 36\\
  \bottomrule
\end{tabular}\\[1mm]
Our pin (\SI{31}{\ohm}) keeps $V_{OL}$ below \SI{0.4}{V} down to about \SI{250}{\ohm}; the specification
demands only \SI{3}{mA} at \SI{0.4}{V}, so \SI{967}{\ohm} is the limit for a weak but compliant device.
With $C_b=\SI{400}{pF}$ and \SI{4.7}{k\ohm}: $t_r=\SI{1.59}{\micro s}$, neither mode is met, at
\SI{100}{kHz} the line does not even reach $V_{DD}$ (\SI{3.0}{V}). $R_P$ has to be reduced (Standard-mode:
$\le\SI{2.9}{k\ohm}$), Fast-mode with \SI{400}{pF} needs $R_P\le\SI{885}{\ohm}<R_{P,min}$. This only works if all devices
meet the additional Fast-mode rating of \SI{6}{mA} at $V_{OL}=\SI{0.6}{V}$ ($R_{P,min}=\SI{450}{\ohm}$, e.g.\
\SI{820}{\ohm}), with Fast-mode Plus devices (\SI{20}{mA}), or if the bus is split by a buffer.
\end{solution}

\task{S/D} Set the input of device~2 to \textsf{L}. What happens to the clock of device~1? Which I$^2$C feature
uses this behaviour? What would happen with two push-pull outputs (part 2)?

\begin{solution}
The bus stays low: the line is the AND of all outputs (``wired AND''). Device~1 can detect this by reading back
the line. I$^2$C uses it for clock stretching (a slow slave holds SCL low), for arbitration between masters (a
master that sends 1 but reads 0 has lost) and for the acknowledge bit. With push-pull outputs the same situation
is a short circuit (part 2).
\end{solution}

\task{D} While the line is low, $R_P$ dissipates $P=V_{DD}^2/R_P$ (\SI{2.3}{mW} at \SI{4.7}{k\ohm}, \SI{11}{mW} at
\SI{1}{k\ohm}). Since $t_r=0.8473\,R_PC_b$, the product $P\cdot t_r=0.8473\,V_{DD}^2C_b$ does not depend on $R_P$.
What follows for the trade-off between speed and power of an open-drain bus, compared to push-pull? Propose a
circuit that speeds up the rising edge without reducing $R_P$, and try it in the simulation (e.g.\ a current
source pull-up or a switched pull-up).

\begin{solution}
$P=V_{DD}^2/R_P$ while low, $t_r=0.8473R_PC_b$, so $P\cdot t_r=0.8473\,V_{DD}^2C_b$ is independent of $R_P$:
halving the rise time doubles the static power. (\SI{4.7}{k\ohm}: \SI{2.3}{mW}, \SI{1}{k\ohm}: \SI{11}{mW}.)
Push-pull has no static current and drives both edges actively. Faster rising edges: a constant-current pull-up
(linear instead of exponential edge, $t_r=0.4\,V_{DD}C_b/I$), a rise-time accelerator that switches on a strong
pull-up as soon as the line rises above a threshold (e.g.\ LTC4311), or a push-pull ``high-speed mode'' as used
by I$^2$C Hs-mode masters. All of them must not fight against a device that pulls low.
\end{solution}

%======================================================================
\needspace{0.45\textheight}
\section{Maximum output voltage: open drain with a 5\,V pull-up}\label{sec:overvoltage}
%======================================================================
\circuitlink{\linkE}{E\_open\_drain\_overvoltage.txt}

A \SI{3.3}{V} microcontroller drives a \SI{5}{V} logic input through an open-drain output and a pull-up to
\SI{5}{V}. PIN1 is a normal GPIO in open-drain mode: the PMOS is switched off, but it is still there. PIN2 is a
\SI{5}{V}-tolerant open-drain pin without a PMOS. The \SI{3.3}{V} regulator can only source current (diode).
The resistor ``MCU load'' represents the current consumption of the chip.

\schematic[0.7]{E}

\task{P/S} Table~5.1 gives the expected high levels and the current that flows from the \SI{5}{V} supply into the
\SI{3.3}{V} supply via PIN1 ($R_{PU}=\SI{4.7}{k\ohm}$, diode forward voltage about \SI{0.7}{V}). Explain the
values, then measure them.

\begin{results}{Table 5.1: open-drain pins with a pull-up to \SI{5}{V} ($R_{PU}=\SI{4.7}{k\ohm}$)}{lccc}
  & high level PIN1 & high level PIN2 & injection current\\ \midrule
  model & $V_{DD}+V_F\approx\SI{4.0}{V}$ & \SI{5.0}{V} & $\dfrac{5-4.0}{4.7}\,\si{mA}\approx\SI{0.2}{mA}$\\
  simulation & \simv{4.10} & \simv{5.00} & \simv{0.19}\\
\end{results}

\begin{solution}
PIN2 rises to \SI{5}{V}. PIN1 is clamped by the body diode of the PMOS to $V_{DD}+V_F\approx\SI{4.0}{V}$; the
current $(5-4.0)\,\si{V}/\SI{4.7}{k\ohm}\approx\SI{0.2}{mA}$ flows into the \SI{3.3}{V} supply.
\end{solution}

\task{D} Is the high level of PIN1 sufficient for a \SI{5}{V} CMOS input with $V_{IH}=0.7\cdot\SI{5}{V}$? Is it
allowed?

\begin{solution}
PIN1: \SI{4.10}{V}, PIN2: \SI{5.00}{V}, injection current \SI{0.19}{mA}. $V_{IH}=\SI{3.5}{V}$: 4.1\,V is
sufficient here, but the pin is outside its absolute maximum rating ($IOVDD+\SI{0.5}{V}$), and the margin shrinks
with a smaller injection current or a hot diode.
\end{solution}

\task{S/D} Set $R_{PU}=\SI{1}{k\ohm}$ and keep the MCU load at \SI{33}{k\ohm} (\SI{100}{\micro A}, e.g.\ in a
sleep mode). Observe $V_{DD}$ (simulation: \underline{\hspace{2cm}}\simv{3.9\,V}) and explain the effect
(\emph{back-powering}): the regulator can only source current, so $V_{DD}$ rises until the chip consumes the
injected current. What happens if the microcontroller is switched off ($V_{DD}=0$) but the \SI{5}{V} side stays on?

\begin{solution}
$V_{DD}$ rises to \SI{3.9}{V} (with \SI{470}{\ohm}: \SI{4.0}{V}), above the absolute maximum of the chip.
While PIN1 is high, \SI{0.25}{mA} flow through the body diode into the \SI{3.3}{V} rail (on average
\SI{0.12}{mA}). The regulator can only source current, so the injected current must be consumed by the chip.
If it consumes less, $V_{DD}$ rises until the load current equals the injected current. With a load of
\SI{1}{mA} (\SI{3.3}{k\ohm}) the rail stays at \SI{3.27}{V}. If the microcontroller is switched off, it is
powered through its pin (``phantom powering''): it may start partially, hold the bus low, or come up in an
undefined state when the supply returns. Data sheets therefore forbid voltages above $V_{DD}+\SI{0.3}{V}$
\ldots$+\SI{0.5}{V}$ on non-tolerant pins, also when the chip is off.
\end{solution}

\task{D} Look up in the data sheets of an RP2040, an RP2350 and an STM32F4 which pins are \SI{5}{V} tolerant
and how large the injection current per pin and in total may be. Name two circuits that shift a
\SI{3.3}{V} open-drain signal to \SI{5}{V} correctly, one of them bidirectional (for I$^2$C).

\begin{solution}
RP2040: no \SI{5}{V} tolerant GPIOs, pin voltage at most $IOVDD+\SI{0.5}{V}$. RP2350: GPIOs are \SI{5}{V}
tolerant while IOVDD is powered (and fail-safe when unpowered), see its data sheet for the exact conditions.
STM32F4: pins marked ``FT'' are \SI{5}{V} tolerant (in input or open-drain mode); the injection current is
limited to about $\pm\SI{5}{mA}$ per pin and $\pm\SI{25}{mA}$ in total, and positive injection is not allowed
on FT pins. Correct level shifters: an FT pin in open-drain mode with a pull-up to \SI{5}{V}; a single
N-channel MOSFET (e.g.\ BSS138) with gate at \SI{3.3}{V} and pull-ups on both sides, which works in both
directions for I$^2$C; dedicated translators (e.g.\ PCA9306, TXS0102).
\end{solution}

%======================================================================
\needspace{0.45\textheight}
\section{Open collector: high voltage, high current, inductive loads}\label{sec:oc}
%======================================================================
\circuitlink{\linkF}{F\_open\_collector\_inductive.txt}

A \SI{3.3}{V} GPIO switches an NPN transistor ($\beta=100$) through the base resistor $R_B$. The collector is an
open-collector output for a \SI{12}{V} load: a relay coil (\SI{0.1}{H}, \SI{400}{\ohm}, \SI{1}{nF} winding
capacitance) or a \SI{12}{V}/\SI{5}{W} lamp. The switch \emph{flyback} adds a free-wheeling diode, opening the
switch \emph{bypass} puts a \SI{24}{V} Zener diode in series with it. The timescale is shortened (switching at
\SI{200}{Hz}) so that you can see several cycles.

\schematic[0.62]{F}

\task{P} Table~6.1 gives the operating point with $R_B=\SI{1}{k\ohm}$. Is the transistor saturated?

\begin{results}{Table 6.1: operating point, relay coil, $R_B=\SI{1}{k\ohm}$}{lccccc}
  & $I_C=\dfrac{\SI{12}{V}}{\SI{400}{\ohm}}$ & $I_B=\dfrac{3.3-0.7}{R_B}$ & $\beta I_B$ & overdrive $\dfrac{\beta I_B}{I_C}$ & $V_{CE}$\\ \midrule
  model & \SI{30}{mA} & \SI{2.6}{mA} & \SI{260}{mA} & 8.7 & $<\SI{0.2}{V}$\\
  simulation & \simv{29.8} & \simv{2.61} & & & \simv{0.07}\\
\end{results}

\begin{solution}
$I_C=\SI{12}{V}/\SI{400}{\ohm}=\SI{30}{mA}$, $I_B=(3.3-0.7)\,\si{V}/\SI{1}{k\ohm}=\SI{2.6}{mA}$,
$\beta I_B=\SI{260}{mA}$, overdrive factor 8.7: the transistor is saturated.
\end{solution}

\task{S/D} Without flyback diode: when the transistor switches off, the coil current continues and charges the
winding capacitance. Energy conservation ($\tfrac12LI^2=\tfrac12CV^2$) gives a peak of
$\SI{12}{V}+I\sqrt{L/C}=\SI{12}{V}+\SI{300}{V}$. Measure it (simulation: \underline{\hspace{2cm}}\simv{296\,V}) and compare with
$V_{CEO}$ of a BC547 (\SI{45}{V}) or of a ULN2003 driver (\SI{50}{V}). What does the simulation not show?

\begin{solution}
$\tfrac12LI^2=\tfrac12CV^2\Rightarrow \Delta V=I\sqrt{L/C}=\SI{30}{mA}\cdot\sqrt{\SI{0.1}{H}/\SI{1}{nF}}=\SI{300}{V}$
above the \SI{12}{V} supply. Simulation: $V_C=\SI{296}{V}$ peak (some energy is lost in the coil resistance),
followed by a ringing that drives the collector to \SI{-14}{V}, where the base-collector diode conducts. This
is far above $V_{CEO}$ of any small-signal transistor. In reality the transistor breaks down (avalanche) at
about \SIrange{45}{80}{V} and absorbs the coil energy in a small volume, which degrades or destroys it; the
simulation models neither breakdown nor the resulting heating, and the real winding capacitance and damping
determine the actual value. The fast edge also radiates (EMI) and couples into the GPIO via the base.
\end{solution}

\task{S} Close \emph{flyback} and measure the maximum collector voltage and the time until the coil current has
fallen to \SI{10}{\percent} (right-click the scope, \emph{Properties}, to change the time base). Then open
\emph{bypass} (diode + \SI{24}{V} Zener) and measure again. Explain why the Zener variant is faster. Why does that
matter for a relay?

\begin{results}{Table 6.2: switching off the relay coil ($\tau=L/R=\SI{0.25}{ms}$)}{lcccc}
  & \multicolumn{2}{c}{diode} & \multicolumn{2}{c}{diode + \SI{24}{V} Zener}\\
  & $V_{C,max}$ & $t_{10\%}$ & $V_{C,max}$ & $t_{10\%}$\\ \midrule
  model & $12+0.7=\SI{12.7}{V}$ & \SI{475}{\micro s} & $12+24+0.7=\SI{36.7}{V}$ & \SI{87}{\micro s}\\
  simulation & \simv{12.8} & \simv{474} & \simv{36.8} & \simv{87}\\
\end{results}
{\small Model: $i(t)=-V_{cl}/R+(I_0+V_{cl}/R)\,e^{-t/\tau}$ with the clamp voltage $V_{cl}=V_F$ or $V_Z+V_F$.}

\begin{solution}
Diode: $V_{C,max}=\SI{12.8}{V}$. The coil current flows on through the diode and decays with
$\tau=L/R=\SI{0.25}{ms}$: $i(t)=-V_F/R+(I_0+V_F/R)\,e^{-t/\tau}$ gives \SI{475}{\micro s} to \SI{10}{\percent}
(simulation: \SI{474}{\micro s}; without $V_F$: $\tau\ln 10=\SI{576}{\micro s}$).\\
Diode + Zener: $V_{C,max}=12+24+0.7\approx\SI{36.8}{V}$ (simulation \SI{36.8}{V}), still below $V_{CEO}$. The
coil now discharges against $V_Z+V_F=\SI{24.7}{V}$:
$i(t)=-\SI{61.8}{mA}+(\SI{30}{mA}+\SI{61.8}{mA})\,e^{-t/\tau}$ reaches \SI{3}{mA} after
$\tau\ln(91.8/64.8)=\SI{87}{\micro s}$ (simulation \SI{87}{\micro s}). The energy is dissipated faster because the
voltage across the coil is higher ($P=V\cdot i$). For a relay the release time determines how fast the contacts
open: a slow decay lets the contacts drift apart slowly, which increases arcing and contact wear. The
designer trades transistor voltage rating against release speed.
\end{solution}

\task{S/D} Increase $R_B$ to \SI{47}{k\ohm}: $I_B=\SI{55}{\micro A}$, so $I_C\le\beta I_B=\SI{5.5}{mA}$ and
$V_{CE}\approx 12-5.5\cdot0.4=\SI{9.8}{V}$. Measure $I_C$ and $V_{CE}$
(simulation: \underline{\hspace{1.5cm}}\simv{5.7\,mA} and \underline{\hspace{1.5cm}}\simv{9.7\,V}). Will the relay still pull in?

\begin{solution}
$I_B=\SI{55}{\micro A}$, $\beta I_B=\SI{5.5}{mA}$: the simulation shows $I_C=\SI{5.7}{mA}$ and
$V_{CE}=\SI{9.7}{V}$. The transistor is in the active region (dissipating \SI{55}{mW}) and the coil receives
only \SI{19}{\percent} of its nominal current: a relay with a typical pull-in current of \SIrange{70}{75}{\percent}
of nominal will not switch. Rule: $R_B$ so that $\beta_{min}I_B\ge 2\ldots 5\cdot I_C$ (overdrive), using the
minimum $\beta$ of the data sheet, not the typical one.
\end{solution}

\task{S/D} Switch the load to the lamp (\SI{12}{V}/\SI{5}{W}, needs \SI{414}{mA}) and complete Table~6.3. A
small-signal transistor in SOT-23 tolerates about \SI{250}{mW}. Why is \SI{330}{\ohm} still not a good solution?
Propose a better driver.

\begin{results}{Table 6.3: lamp load}{lcccc}
  & $I_B$ & $I_C$ & $V_{CE}$ & $P=V_{CE}I_C$\\ \midrule
  $R_B=\SI{1}{k\ohm}$, model & \SI{2.6}{mA} & $\beta I_B=\SI{260}{mA}$ & $12-0.26\cdot29=\SI{4.5}{V}$ & \SI{1.2}{W}\\
  $R_B=\SI{1}{k\ohm}$, simulation & \simv{2.56} & \simv{256} & \simv{4.57} & \simv{1.17}\\
  $R_B=\SI{330}{\ohm}$, model & \SI{7.9}{mA} & \SI{414}{mA} (saturated) & $\approx\SI{0.1}{V}$ & $\approx\SI{50}{mW}$\\
  $R_B=\SI{330}{\ohm}$, simulation & \simv{7.72} & \simv{410} & \simv{0.12} & \simv{0.05}\\
\end{results}

\begin{solution}
$R_B=\SI{1}{k\ohm}$: $I_C=\beta I_B=\SI{256}{mA}$ instead of \SI{414}{mA}, $V_{CE}=\SI{4.6}{V}$,
$P=\SI{1.17}{W}$: the lamp stays dim and the transistor overheats within seconds.
$R_B=\SI{330}{\ohm}$: $I_B=\SI{7.7}{mA}$, saturated, $I_C=\SI{410}{mA}$, $V_{CE}=\SI{0.12}{V}$, $P=\SI{50}{mW}$.
But \SI{7.7}{mA} is a substantial load for the GPIO ($V_{OH}$ drops, part 1), the overdrive is only 1.9 with
the typical $\beta$ and becomes insufficient with the minimum $\beta$ at high current, and the cold filament
draws about ten times the nominal current when switched on. Better: a logic-level N-MOSFET (e.g.\ with
$R_{DS(on)}<\SI{0.1}{\ohm}$ at $V_{GS}=\SI{2.5}{V}$, gate resistor about \SI{100}{\ohm}, pull-down
\SI{100}{k\ohm}), a Darlington array such as the ULN2003 (\SI{500}{mA}, integrated clamp diodes), or a
high-side switch with current limit.
\end{solution}

\task{D} Summarise: what limits the maximum output voltage and the maximum output current of a push-pull pin, of
an open-drain pin and of an external open-collector stage?

\begin{solution}
Push-pull: $V_{OH}\le V_{DD}$ (body diodes clamp at $V_{DD}+V_F$), current limited by $R_{on}$, $I_{sat}$ and the
per-pin/per-domain ratings (part 1). Open-drain pin: voltage limited by the pin structure (body/ESD diode to
$V_{DD}$, or the rating of an FT pin, typically \SI{5.5}{V}), sink current by the NMOS like push-pull; the high
level is only as strong as the pull-up. External open collector/drain: voltage limited by $V_{CEO}$/$V_{DS,max}$ of
the transistor (and inductive kick-back), current by $\beta I_B$ (or $V_{GS}$ for a MOSFET), the
saturation voltage and the power dissipation of the package.
\end{solution}

%======================================================================
\needspace{0.45\textheight}
\section{Fast edges on a PCB trace}
%======================================================================
\circuitlink{\linkG}{G\_transmission\_line.txt}

The push-pull pin drives a \SI{50}{\ohm} trace with \SI{1}{ns} delay (about \SI{15}{cm} on FR4). At the end there is
the input of another chip: \SI{5}{pF} and ESD diodes to $V_{DD}$ and GND. $R_S$ is a series resistor at the
driver. The output stage is set to its highest drive strength and fastest slew rate, modelled by
$\beta_N=\SI{30}{mA/V^2}$ and $\beta_P=\SI{25}{mA/V^2}$. The slew rate is adjustable as on many
microcontrollers (RP2040: \texttt{SLEWFAST}, STM32: \texttt{OSPEEDR}).

\schematic[0.75]{G}

\task{P} A trace is ``electrically long'' if the edge is shorter than about twice its delay. Table~7.1 gives the
values for $R_S=0$ and the fastest slew rate. Explain the sequence: launched step, doubling at the open end,
clamping, reflection at the driver.

\begin{results}{Table 7.1: reflections for $R_S=0$ ($R_{on,P}=1/(\SI{25}{mA/V^2}\cdot\SI{2.7}{V})=\SI{15}{\ohm}$)}{lc}
  edge time / round trip & $<\SI{1}{ns}$ / \SI{2}{ns}: electrically long\\
  $\Gamma_S=(R_{on}+R_S-Z_0)/(R_{on}+R_S+Z_0)$ & $(15-50)/(15+50)=-0.54$\\
  $\Gamma_L$ (open end, \SI{5}{pF} only rounds the edge) & $\approx+1$\\
  launched step $V_{DD}\,Z_0/(Z_0+R_{on})$ & \SI{2.54}{V}\\
  first voltage at RX (before clamping) & $2\cdot\SI{2.54}{V}=\SI{5.1}{V}$\\
  matched series resistor $R_S=Z_0-R_{on}$ & $\approx\SI{36}{\ohm}$ (\SI{33}{\ohm} standard value)\\
\end{results}

\begin{solution}
In its fastest setting the output edge takes well below \SI{1}{ns}, the round trip takes
\SI{2}{ns}: the trace is electrically long. With $R_{on,P}=1/(\SI{25}{mA/V^2}\cdot\SI{2.7}{V})=\SI{15}{\ohm}$:
$\Gamma_S=(15-50)/(15+50)=-0.54$, $\Gamma_L\approx+1$ (open end; the \SI{5}{pF} only rounds the edge).
Launched step $\SI{3.3}{V}\cdot 50/(50+15)=\SI{2.54}{V}$; at the receiver it doubles to \SI{5.1}{V}, i.e.\
\SI{1.8}{V} above $V_{DD}$, so the ESD diode clamps it. The reflected wave returns to the driver, is reflected
with $-0.54$ and arrives as an undershoot of about $-0.54\cdot\SI{2.5}{V}$ after another \SI{2}{ns}
(again limited by the GND clamp): the edge rings with a period of \SI{4}{ns}.
\end{solution}

\task{S} Measure overshoot and undershoot at RX and the peak currents in the ESD diodes for the values of $R_S$ in
Table~7.2. Which $R_S$ gives the cleanest edge? Compare with $R_S=Z_0-R_{on}$.

\begin{results}{Table 7.2: series termination (fastest slew rate)}{lccccc}
  $R_S$ & \SI{0}{\ohm} & \SI{10}{\ohm} & \SI{22}{\ohm} & \SI{36}{\ohm} & \SI{50}{\ohm}\\ \midrule
  $V_{RX,max}$, simulation & \simv{3.92} & \simv{3.89} & \simv{3.74} & \simv{3.30} & \simv{3.30}\\
  $V_{RX,min}$, simulation & \simv{-0.64} & \simv{-0.61} & \simv{-0.52} & \simv{-0.02} & \simv{0.00}\\
  $I_{ESD}$ to $V_{DD}$ / GND, simulation & \simv{30/40\,mA} & \simv{14/21} & \simv{1/4} & \simv{0} & \simv{0}\\
\end{results}

\begin{solution}
$R_S=0$: RX swings to \SI{3.9}{V} and \SI{-0.64}{V}, i.e.\ beyond the \SI{\pm0.5}{V} limits; the ESD
diodes conduct up to \SI{30}{mA} (to $V_{DD}$) and \SI{40}{mA} (to GND) for each edge, and the ringing can
produce double clocking at the receiver.
\begin{tabular}{lccccc}
  \toprule
  $R_S$ / \si{\ohm} & 0 & 10 & 22 & 36 & 50\\
  \midrule
  $V_{RX,max}$ / V & 3.92 & 3.89 & 3.74 & 3.30 & 3.30\\
  $V_{RX,min}$ / V & $-0.64$ & $-0.61$ & $-0.52$ & $-0.02$ & 0.00\\
  $I_{ESD}$ / mA & 30/40 & 14/21 & 1/4 & 0 & 0\\
  \bottomrule
\end{tabular}\\[1mm]
The cleanest edge is at $R_S\approx\SI{36}{\ohm}=Z_0-R_{on}$ (between $50-15$ for the rising and $50-12$ for the
falling edge; standard value \SI{33}{\ohm}). Then the driver launches $V_{DD}/2$, the open end doubles it to
$V_{DD}$, and the reflection is absorbed at the source. Larger $R_S$ only slows the edge (the line voltage
steps up in several round trips).
\end{solution}

\task{S/D} Set $R_S$ back to 0 and reduce the slew rate. From which slew rate on does the ringing disappear?
Why should you select the slowest slew rate that meets the timing of your interface?

\begin{solution}
The slider sets the slew rate of the gate; the output switches during only part of the gate swing, so
measure the edge at PIN. At \SI{1}{V/ns} the ringing is unchanged (ESD currents \SI{19}{mA}/\SI{24}{mA}),
at \SI{0.3}{V/ns} it is reduced (\SI{4.5}{mA}/\SI{6.6}{mA}), at \SI{0.1}{V/ns} (output edge several
\si{ns}, much longer than the \SI{2}{ns} round trip) the trace behaves like a lumped capacitance and the edge
is clean. A slow edge also has a narrower spectrum: less crosstalk, less radiated EMI, smaller ground bounce
when many pins switch together, and less supply current (part 3). The fastest setting is only needed when the
timing budget of the interface requires it.
\end{solution}

%======================================================================
\needspace{0.45\textheight}
\section{Design task}
%======================================================================

\task{P/S/D} A \SI{3.3}{V} microcontroller board has to
\begin{enumerate}[label=(\alph*)]
  \item read a \SI{5}{V} sensor over I$^2$C at \SI{400}{kHz}; the cable is \SI{30}{cm} long, the total bus
        capacitance is estimated at \SI{250}{pF},
  \item switch a \SI{24}{V}/\SI{50}{mA} solenoid valve at up to \SI{10}{Hz},
  \item send a \SI{10}{MHz} SPI clock over a \SI{10}{cm} trace to a flash memory.
\end{enumerate}
For each signal choose the output type and configuration of the pin (push-pull/open drain, drive strength, slew
rate), dimension the external components and verify your choice in a simulation that you build from the given
circuits. Document the remaining risks.

\begin{solution}
A possible solution (other justified solutions are equally valid):
\begin{enumerate}[label=(\alph*)]
  \item I$^2$C: open drain is mandatory (wired AND, clock stretching, ACK). The RP2040 pins are not
        \SI{5}{V} tolerant $\Rightarrow$ a BSS138 level shifter per line (or a PCA9306) with pull-ups on both
        sides. The cable (\SI{250}{pF}) is on the \SI{5}{V} side: $R_{P,max}=\SI{300}{ns}/(0.8473\cdot\SI{250}{pF})
        =\SI{1.42}{k\ohm}$, but $R_{P,min}=(\SI{5}{V}-\SI{0.4}{V})/\SI{3}{mA}=\SI{1.53}{k\ohm}$: no resistor meets
        both. Options: devices that are rated for \SI{6}{mA} at \SI{0.6}{V} ($R_{P,min}=\SI{733}{\ohm}$, choose
        \SI{1.2}{k\ohm}), a rise-time accelerator or bus buffer, or Standard-mode at \SI{100}{kHz}
        ($R_P=\SI{1.8}{k\ohm}\ldots\SI{4.7}{k\ohm}$). On the short \SI{3.3}{V} side \SI{4.7}{k\ohm} is enough.
        Simulation: part D with \SI{5}{V} pull-up and \SI{250}{pF}. Risk: the capacitance estimate; measure
        $t_r$ on the real board.
  \item Solenoid: low-side N-MOSFET (logic level, $V_{DS}\ge\SI{40}{V}$) or ULN2003 channel; GPIO push-pull,
        \SI{100}{\ohm} gate resistor, \SI{100}{k\ohm} pull-down so the valve stays off during reset (pins are
        inputs then). Free-wheeling: diode + Zener (e.g.\ \SI{24}{V}) if the valve must close quickly,
        $V_{DS}\ge 24+24+1\approx\SI{50}{V}$, otherwise a diode alone. Simulation: part F with \SI{24}{V},
        \SI{480}{\ohm} and the valve inductance.
  \item SPI clock: push-pull, moderate drive strength and slew rate, series resistor $R_S=Z_0-R_{on}$
        (e.g.\ \SI{33}{\ohm}) at the MCU. A \SI{10}{cm} trace has \SI{0.67}{ns} delay; with the fastest edge
        it is long, so terminate or slow the edge (still $t_r\ll\SI{50}{ns}$ half period). Simulation: part G with
        the trace delay changed to \SI{0.67}{ns}. Risk: ringing causes double clocking of the flash.
\end{enumerate}
\end{solution}

\end{document}
